\documentclass[11pt,a4paper]{article}
\usepackage[margin=27mm,headheight=15pt]{geometry}
\usepackage{fontspec}
\setmainfont{Latin Modern Roman}
\setsansfont{Latin Modern Sans}
\setmonofont{Latin Modern Mono}[Scale=MatchLowercase]
\usepackage{amsmath,amssymb,amsthm,mathtools,mathrsfs}
\usepackage{microtype,graphicx,booktabs,tabularx,longtable,array,enumitem}
\usepackage[dvipsnames]{xcolor}
\usepackage[round,authoryear]{natbib}
\usepackage{hyperref}
\usepackage{bookmark}
\hypersetup{colorlinks=true,linkcolor=MidnightBlue,citecolor=MidnightBlue,urlcolor=MidnightBlue,pdftitle={Split-zero structures and the zeta-function programme},pdfauthor={KokunoYumeto programme; AI-assisted editorial companion}}
\setlength{\emergencystretch}{2em}
\setlength{\parskip}{.35em}
\setlist{itemsep=.3em,topsep=.4em}
\allowdisplaybreaks[2]
\newtheorem{theorem}{Theorem}[section]
\newtheorem{proposition}[theorem]{Proposition}
\newtheorem{lemma}[theorem]{Lemma}
\newtheorem{corollary}[theorem]{Corollary}
\theoremstyle{definition}
\newtheorem{definition}[theorem]{Definition}
\newtheorem{example}[theorem]{Example}
\theoremstyle{remark}
\newtheorem{remark}[theorem]{Remark}
\title{Split-zero structures and the zeta-function programme\\[.5em]\large A mathematical exposition with human-source provenance}
\author{An AI-assisted companion to the KokunoYumeto / The Clankers programme}
\date{17 September 2026\\\small Fixed mathematical baseline: the accepted 726-page edition}
\begin{document}
\maketitle
\begin{abstract}
This companion explains a research programme connecting support-sensitive scalar structures, a classical theta quotient, complete zero jets, exponential periods, and determinant estimates. It defines the maps that join these constructions and identifies their human mathematical sources. The programme's theta quotient is explicitly identified with Ralf Meyer's quotient; its relationship to the geometric results of Pierre Deligne is stated with the relevant hypotheses. The current determinant calculation and its remaining target are written on their original spaces. The companion retains the complete proof archive and adds an entry route, source attributions, and a record of what has and has not been checked. The programme remains open; no proof of the Riemann hypothesis is claimed.
\end{abstract}
\clearpage
\section*{Purpose and reading conventions}
Human readability and accurate provenance are commitments of this
programme. Its mathematics should be pedagogically articulated so that people
can understand and examine it. Its source record should fairly credit the
humans whose work it uses. The Leiden Declaration and the September
Math and AI declaration provide context for these editorial concerns
\citep{AlperEtAl2026Leiden,AvilaEtAl2026Misalignment}.
Neither declaration supplies mathematical hypotheses or endorses this programme.

The intended reader knows linear algebra, basic complex analysis, and the
language of vector spaces and quotients. The analytic and categorical parts
introduce their objects before using them. The appendix spells out the
support, complex, and source-lifting maps that connect the opening definitions
to the later finite calculations. A reader chiefly interested in the current
estimate can follow the main text's period and metric constructions, then
return to that appendix for the analytic realization.

For stable historical page references, \textbf{R14} refers to the complete editable source
\nolinkurl{14_CURRENT_KERNEL_AND_MULTIPLICITY_PROOFS.tex};
\textbf{R09} refers to \nolinkurl{09_UPDATED_JOINT_NOTE.tex} in the same
accepted edition. Page numbers attached to R14 refer to the printed pages of
the preserved 692-page predecessor reader, \emph{Original Mixed Cohomology: Periods, Kernels and
Arithmetic Control}. The version is available with complete public sources in
the \href{https://github.com/KokunoYumeto/zeta-function-research-reader/releases/tag/rgc-growing-degree-2026-09-17-de064490}{17 September growing-degree release}.
Its source commit is \texttt{de0644902496e1718879f4b625f9ff9faafc5ae5}.
These fixed predecessor locators remain meaningful. The accepted successor,
\textbf{R726}, contains 726 pages and preserves every predecessor source byte
while inserting the new proofs and their earlier receiving calculations.
Its source commit is \texttt{be79088cb1fa4d79b31724856eefcdbd0839b8cc};
its publication is \href{https://doi.org/10.5281/zenodo.22803503}{DOI 10.5281/zenodo.22803503}.
Section~\ref{sec:moving-cutoff} explains its advance with R726 locators.

Results imported from that archive are identified by their proof locations.
Elementary comparison statements established here carry their proofs.
References to the literature identify what is inherited, with conventions and
hypotheses wherever those affect the map. A proved project calculation is not
thereby a claim of historical novelty. Exact source-verification levels and
unresolved bibliographic trails are retained in the accompanying human-source
guide and attribution records.

The originating programme is directed under the public name KokunoYumeto.
Its notes and this companion use AI assistance. The accompanying disclosure
records that role separately from the authors of the mathematical literature.
No independent human peer review of this companion is asserted. The reader
should judge the definitions, proofs, source comparisons, and scoped evidence
provided here and in the complete archive.

\clearpage
{\small\setlength{\parskip}{0pt}\tableofcontents}
\clearpage
% Additive human-facing mathematical route; fixed accepted 692-page baseline.
\section{The question and the present answer}\label{the-question-and-the-present-answer}

The programme asks whether a cohomological construction over a proposed absolute base can retain the arithmetic information of the Riemann zeta function and provide sufficiently strong control of a specified spectral action. The current calculations study the consequences of a hypothetical off-critical-line zero quartet. They build the associated spaces and maps, preserve repeated zeros, and calculate norms and determinant returns on those actual spaces. This is research toward a geometric argument about zeta zeros. The edition proves a substantial collection of finite-dimensional identities and estimates. It does not prove the Riemann hypothesis, establish the existence of an off-line quartet, or exclude such a quartet.

The inspiration from Pierre Deligne is precise at the level of ambition: a cohomological realization should come with a structure that controls spectral size. Deligne's theorem for smooth projective varieties over finite fields has an actual variety, an actual Frobenius action and a purity statement. Here the base, support structure, analytic classes, periods and metric comparisons have to be constructed and connected. Their formal resemblance to finite-field geometry is a reason to calculate those connections, not a proof of a purity theorem. The original audit constructs specific scaling and nilpotent maps and records the differences between them. Deligne's \emph{La conjecture de Weil. I}, theorem 1.6, supplies that finite-field statement \cite{Deligne1974WeilI}. His later theory of weights is a distinct development \cite{Deligne1980WeilII}. Neither statement is silently imported into the present analytic construction.

The central allocation problem, retained in the 726-page successor, is concrete. A known total determinant growth splits into two correlated contributions. Their sum has been evaluated. Their separate leading allocation has not. The surviving expression is a signed spectral profile of eight specified four-by-four positive matrices, together with four specified complementary quotient returns. The matrices are computed from the original period, Gamma recurrence, conductor division and lower-evaluation minimum; they are not arbitrary test matrices. A later part of this exposition writes down their exact formula.

There is a second, related calculation concerning the original action's canonical allowance. Strong lower bounds for that allowance have now been proved, including growing polynomial-degree windows. An improvement to the mixed allocation does not automatically improve that action bound: an exact cancellation must be applied first. Keeping these two questions connected through that identity is essential to understanding both what has been achieved and what remains.

\section{What the support base is recording}\label{what-the-support-base-is-recording}

The base notation distinguishes \emph{absence} from a \emph{present coefficient equal to zero}. This project construction is documented in \cite{GlobalizationNote2026,SplitZeroSecondary2026}; its human mathematical background and exact maps are developed in Appendix~\ref{app:scalar-and-base}. For a ring \(R\), the underlying carrier is

\[
G(R)=\{\tau\}\sqcup R^\bullet,\qquad e_R=0_R^\bullet.
\]

The two observations are

\[
(p_R,b_R)(\tau)=(0,0),\qquad
(p_R,b_R)(r^\bullet)=(r,1).
\]

Their joint map is injective. Indeed two present points with the same first coordinate have the same coefficient; an absent and a present point have different second coordinates. Thus \(p_R(e_R)=p_R(\tau)=0\) does not identify \(e_R\) with \(\tau\). The support bit records a distinction that the scalar observation forgets. The proposed absolute base is denoted \(\mathbf F_{1,\tau}=\{\tau,1\}\), with its specified pointed/preaddition structure. A structural map sends its absent element to \(\tau\) and its unit to \(1_R^\bullet\). The structural base map belongs to the pointed/preaddition category; the coefficient adjunction belongs to the diagram category. Their complete definitions and proofs are retained in the original base-formation source.

Why does this distinction survive so far into an analytic calculation? Because a relation may map to zero \emph{at an existing support label}. For a coefficient map \(L\), its lift on a present label \(\omega\) is

\[
(\omega,x)\longmapsto(\omega,Lx).
\]

If \(Lx=0\), the image is \((\omega,0)\), not external absence. When a larger coefficient face is obtained by zero insertion \(j_A^B\), the actual compatibility required is \(L^B j_A^B=j_A^B L^A\). This is a commutative-square statement about the maps. It preserves the original outer leg and every coefficient face; it does not promote a coefficient-linear section to a unital algebra morphism. {[}R09, §1, ``Scope, original input, and retained support''; R14, AKP1--5, §26, pp.110--111. The full categorical definitions and proofs are in the original base-formation audit, chapter 01, §§1--4.{]}

\section{The analytic source: theta, Mellin transform and full jets}\label{the-analytic-source-theta-mellin-transform-and-full-jets}

Write the Riemann xi function in the convention

\[
\xi(s)=\frac12s(s-1)\pi^{-s/2}\Gamma(s/2)\zeta(s),
\qquad g(s)=2\xi(s).
\]

The Riemann zeta and theta setting is classical \cite{Riemann1859}. The original test space and theta map are

\[
V=\{\phi\in\mathcal S(\mathbb R)_{\mathrm{ev}}:
          \phi(0)=0,\ \int_{\mathbb R}\phi(x)\,dx=0\},
\qquad \Theta\phi(x)=2\sum_{n\ge1}\phi(nx).
\]

The target \(\mathscr B\) consists of smooth functions on \(\mathbb R_{>0}\) for which

\[
\sup_{x>0}x^b|D^jF(x)|<\infty
\quad(b\in\mathbb Z,\ j\ge0),\qquad D=-x\partial_x.
\]

Its Mellin transform is \(\mathcal MF(s)=\int_0^\infty F(x)x^s\,dx/x\). Integration by parts, whose two boundary terms vanish by the defining decay at zero and infinity, gives

\[
\mathcal M(DF)(s)=s\mathcal MF(s).
\]

This is why a polynomial in the Euler operator becomes the same polynomial in the arithmetic variable \(s\). The quotient \(Q=\mathscr B/\Theta V\) records analytic classes. The original two-term coefficient resolution is

\[
[V\oplus V\xrightarrow{\ (\phi,\psi)\mapsto
          \Theta(\phi-\widehat\psi)\ }\mathscr B].
\]

Its chart and tensor-resolution identifications are proved in the original audit; the current arithmetic map uses its top tensor cohomology \(Q^{\otimes k}\). {[}R09 §1; R14 AKP1--2.{]}

This is a direct debt to Ralf Meyer's spectral construction, not merely a generic similarity to the literature. In Meyer, \emph{A spectral interpretation for the zeros of the Riemann zeta function}, §3, Definition 3.1 and equation \(11\), \(Zf=\sum_{n\ge1}f(nx)\), the even Schwartz subspace has \(f(0)=\widehat f(0)=0\), and the Poisson identity supplies its two asymptotic conditions. Here \(V\) is that subspace under the stated Fourier convention and \(\Theta=2Z\). In fact the original quotient \(Q\) is topologically and equivariantly Meyer's quotient \(\mathcal H_-^0\), as the explicit comparison below shows. It should not be presented as a newly invented quotient. The source conventions are related as follows: Meyer's completed zeta denoted \(\xi\) does not contain the same \(s(s-1)\) factor as the programme's \(g\), and the temporary generator \(x\partial_x\) in his Proposition 3.2 has the opposite sign from \(D\). His closed-range result is Theorem 3.3. Meyer in turn explicitly discusses Alain Connes's earlier spectral interpretation. These are specific debts to Meyer and the antecedent he identifies in Connes \cite{Meyer2005,Connes1999}. The version used for the convention comparison is Meyer's pinned arXiv version 3, §§2--4.

For completeness, put \(f(t)=F(e^t)\). The source seminorms become \(p_{b,j}(F)=\sup_te^{bt}|f^{(j)}(t)|\). Meyer requires \(e^{st}f(t)\) to be Schwartz for each real \(s\). If \(b_-<s<b_+\) are integers, every polynomial weight \((1+|t|)^n\) is bounded by a constant times \(e^{(b_+-s)t}\) on \(t\ge0\), and by a constant times \(e^{(b_--s)t}\) on \(t\le0\). Leibniz's rule therefore gives

\[
\sup_t(1+|t|)^n|\partial_t^j(e^{st}f)|
\le C\sum_{r=0}^j\binom jr|s|^{j-r}
                  (p_{b_-,r}(F)+p_{b_+,r}(F)).
\]

Conversely \(e^{bt}f^{(j)}=(\partial_t-b)^j(e^{bt}f)\), so each \(p_{b,j}\) is bounded by a finite sum of Meyer's Schwartz seminorms. The identity is consequently a homeomorphism \(\mathscr B\simeq\mathcal H_-\). Since \(2V=V\), its quotient map is the literal identity \(Q\to\mathcal H_-^0,[F]\mapsto[F]\). Both spaces carry \(U_aF(x)=F(x/a)\), so this map also intertwines their dilation representations. This exact precedent determines the credit at the first definition of \(Q\).

Now stipulate a complete zero quartet

\[
\rho=\tfrac12+\delta+i\gamma,\qquad
0<\delta<\tfrac12,\quad\gamma>2,\quad m\ge1,
\] \[
\mathcal R=\{\rho,\bar\rho,1-\rho,1-\bar\rho\},\qquad
h(s)=\prod_{\alpha\in\mathcal R}(s-\alpha)^m.
\]

The hypothesis specifies that \(m\) is the complete order at each member of this packet. Consequently \(v_h=g/h\) is entire at those four points and \(v_h(\alpha)\ne0\). The primary algebra \(E_h=\mathbb C[s]/(h)\) has dimension \(4m\). It retains derivatives through order \(m-1\), not just the four values. Multiplication by

\[
\upsilon_h=j_hv_h\in E_h
\]

is an automorphism because its constant value in each local primary factor is nonzero. Its matrix in raw derivative coordinates is triangular, with entries

\[
\binom rs v_h^{(r-s)}(\alpha)\qquad(r\ge s).
\]

The off-diagonal entries matter whenever \(m>1\). Replacing this matrix by four scalar values would change the map. The original Euler-section proof constructs \(F_h\) with \(\mathcal MF_h=v_h\) and an exact jet section \(\sigma_h:E_h\hookrightarrow Q\). It proves those maps on their stated domains. {[}R09 §1; R14 OCS1, OCS12--18, AKP2; full analytic proof in original audit chapter 04.{]}

\section{Why the sum algebra is smaller than the full tensor product}\label{why-the-sum-algebra-is-smaller-than-the-full-tensor-product}

Fix \(k=4l+1\), \(l\ge1\), and retain

\[
c=k/2,\qquad e_k=1+k(m-1),\qquad
\lambda_{ab}=c+(2a-k)\delta+i(2b-k)\gamma,
\] \[
\chi(S)=\prod_{a,b=0}^{k}(S-\lambda_{ab})^{e_k},\qquad
q=e_k(k+1)^2,\qquad E=\mathbb C[S]/(\chi).
\]

The full ordered tensor product is \(E_h^{\otimes k}\), of dimension \((4m)^k\). The programme does not identify it with the \(q\)-dimensional algebra \(E\). It gives the exact injection

\[
\eta:E\longrightarrow E_h^{\otimes k},\qquad
\eta[P]=\upsilon_h^{\otimes k}P(s_1+\cdots+s_k)1.
\tag{A}
\]

Here is the reason for the exponent \(e_k\), including its constant. In a fixed ordered primary tensor block write \(s_i=\alpha_i+n_i\), where \(n_i^m=0\). The nilpotents commute. Every monomial in \((\sum_i n_i)^{k(m-1)+1}\) has some exponent at least \(m\), so that power vanishes. In the preceding power the coefficient of \(n_1^{m-1}\cdots n_k^{m-1}\) is

\[
\frac{[k(m-1)]!}{[(m-1)!]^k},
\]

which is nonzero over \(\mathbb C\); this is the only surviving monomial of that total degree. The nilpotence index is therefore exactly \(e_k\). The possible sums of the original roots are precisely the distinct \(\lambda_{ab}\). Their real and imaginary parts separately distinguish \(a,b\), since \(\delta,\gamma>0\). The annihilator of the cyclic unit for multiplication by \(\sum_i s_i\) is thus exactly \(\chi\). Polynomial substitution has kernel \((\chi)\), and multiplication by the invertible full unit \(\upsilon_h^{\otimes k}\) preserves injectivity. This proves (A) and the intertwining identity

\[
\eta M_S=\Bigl(\sum_iM_{s_i}\Bigr)\eta.
\]

This calculation illustrates the intended role of an exact map: the smaller cyclic module is a specified submodule of the complete tensor object. It does not erase the other tensor directions. {[}R14 OCS1--3, §1 pp.20--21; CAI10, §198 p.624.{]}

At the function level the corresponding source is

\[
\mathcal V_{h,k}(P)=P(D_1+\cdots+D_k)F_h^{\otimes k},\qquad
q^{(k)}\mathcal V_{h,k}(P)=\sigma_h^{\otimes k}\eta[P].
\]

The relation \(P\mapsto P+\chi\kappa\) changes the representative by an explicit coboundary. Ordered monic division gives polynomials \(Q_a(\mathbf s)\) with

\[
\chi(s_1+\cdots+s_k)=\sum_{a=1}^kh(s_a)Q_a(\mathbf s).
\]

The final remainder vanishes by the annihilator calculation. With the original \(\phi_0\) satisfying \(\Theta\phi_0=h(D)F_h\), placed in the degree-zero plus leg as \((\phi_0,0)\), define

\[
\Psi_\kappa=\sum_{a=1}^k(-1)^{a-1}Q_a(\mathbf D)
\kappa\!\left(\sum_bD_b\right)
\bigl(F_h^{\otimes(a-1)}\otimes\phi_0
                       \otimes F_h^{\otimes(k-a)}\bigr).
\]

In the \(a\)-th summand the tensor differential contributes another sign \((-1)^{a-1}\); the two signs multiply to \(1\). Applying the differential therefore gives exactly \(d\Psi_\kappa=\mathcal V_{h,k}(\chi\kappa)\). This proves representative independence with its actual homotopy. By contrast, lying in a subsequent observation kernel does not make a class a coboundary: the original injection may still carry it to a nonzero arithmetic class. {[}R14 AKP2--5, pp.110--111.{]}

\section{What the period observation sees}\label{what-the-period-observation-sees}

The additional observation uses the original polynomial primitive \(\Phi_h'=h\), \(\Phi_h(0)=0\), a nonzero parameter \(u\), a fixed branch, and the oriented contours

\[
\Gamma_j=-\ell_0+\ell_j,\qquad
\arg\ell_j=\frac{\arg u+\pi+2\pi j}{4m+1}.
\]

On a degree-less-than-\(4m\) representative,

\[
(\Pi[P])_j=\int_{\Gamma_j}e^{\Phi_h(z)/u}P(z)\,dz.
\]

The preceding full period theorem supplies convergence and invertibility on its stated domain. In file 03, PDE1--13, the calculation uses the actual determinant differential equation and its Gamma/Fourier base point; this proof is given directly, rather than invoking an unspecified abstract period theorem. It is not permissible to choose period entries independently: each is this integral of the specified representative. Let \(n=4m\), \(N=n+1\), \(U=M_{\upsilon_h}\) and \(\iota[P]=[P(s_1+\cdots+s_k)]\). The observed map is

\[
L=P_k\Pi^{\otimes k}U^{\otimes k}\iota=\jmath\Lambda,
\quad P_k=\frac1N\sum_{a=0}^{N-1}(C^{-a})^{\otimes k},
\quad K=\ker\Lambda,
\]

where \(\Lambda:E\twoheadrightarrow B\) is onto its actual image and \(\jmath\) includes that image into the ordered period tensor space. The literal cyclic matrix is

\[
(Cx)_j=x_{j+1}-x_1\ (j<N-1),\qquad
(Cx)_{N-1}=-x_1.
\]

With \(\zeta=e^{2\pi i/N}\), \(f_t=(\zeta^{jt}-1)_{j=1}^{N-1}\), one computes \(Cf_t=\zeta^tf_t\). The finite average selects the symmetric tensor words whose sum of indices is zero modulo \(N\). However, these Fourier columns have Gram \(N(I+\mathbf1\mathbf1^*)\), so this coordinate diagonalization is not an assertion that the projector is orthogonal in the original metric. The current source retains the induced full tensor Gram and all its cross terms. {[}R14 OCS2--11, §1 pp.20--23.{]}

For the simple-quartet case \(m=1\), the corresponding period curve is a quadric with an actual cyclic orbit. Its five-member orbit gives

\[
\operatorname{rank}\Lambda=k^2-6k+17,\qquad
\dim K=8k-16.
\]

The actual period calculation supplies a radius \(R_*\) forcing the relevant orbit on its stated branch domain; the metric applications retain \(|u|\ge R_*\), and later pole presentations have their further specified exclusions or sufficient radius. These numbers are not asserted for arbitrary multiplicity. The all-\(m\) geometry retains its own Segre--Veronese maps, nilpotents and stabilizer. {[}R14 OQX3, §3 pp.32--33; PQO/RPC, §§4--5 pp.34--42; OMG, §13 pp.84ff.{]}

Ranks alone do not determine the desired volumes. The latter depend on the period coefficients and on how their kernel sits inside the original source metric. This is the reason the calculation continues after the exact rank theorem.

\section{The arithmetic norm and the two Gamma references}\label{the-arithmetic-norm-and-the-two-gamma-references}

The arithmetic density is

\[
w_h(y)=\frac{|v_h(\tfrac12+iy)|^2}{2\pi},\qquad
m_{h,k}=w_h^{*k}.
\]

For \(P\in\mathbb C[S]\), the finite source norm is

\[
\|P\|_{\mathrm{ar},k}^2
=\int_{\mathbb R}|P(c+iy)|^2m_{h,k}(y)\,dy.
\]

It comes from the same Mellin/theta class; no replacement probability measure is implicit. In particular the full mass is retained. The analytically tractable reference measures are

\[
\sigma(y)=\frac{|\Gamma(\tfrac14+iy/2)|^2}{2\pi},
\quad c_\sigma=\sqrt{2\pi},\quad
c_a=2^{1-2a}\Gamma(2a),
\] \[
r_k^0(y)=\frac{c_\sigma^k}{c_{k/4}}
          \frac{|\Gamma(k/4+iy/2)|^2}{2\pi}.
\]

Their masses are \(c_\sigma\) and \(c_\sigma^k\), respectively. The monic Gamma-polynomial norm convention is \(h_{n,s}=(2\pi)^{s/2}n!(s/2)_n\), with \(s=1\) or \(k\). These are classical Gamma/Meixner--Pollaczek structures, used here with the displayed scaling and mass convention \cite{DLMF18Current,KoekoekLeskySwarttouw2010}.

An exact identity explains why two source orders can be compared on the same arithmetic quotient. Set

\[
b_j=2j+\tfrac12,\qquad
f_l(S)=\prod_{j=0}^{l-1}(S-(c+b_j)),\qquad
C_k=\frac{c_\sigma^k}{c_{k/4}}4^{-l}.
\]

The classical Gamma recurrence \cite{DLMF5Current} gives

\[
\frac{\Gamma(l+1/4+iy/2)}{\Gamma(1/4+iy/2)}
=(i/2)^l\prod_{j=0}^{l-1}(y-ib_j).
\]

Taking modulus on the real line gives the complete equality

\[
r_k^0(y)=C_k|f_l(c+iy)|^2\sigma(y).
\tag{B}
\]

The holomorphic recurrence uses the displayed minus signs before taking the modulus. For example \(k=5\) gives \(f_1(S)=S-3\), \(C_5=16\pi^2/3\). Since \(k\) is odd, each original \(\lambda_{ab}\) has nonzero imaginary part, whereas all roots of \(f_l\) are real. Thus \(f_l\) and \(\chi\) are coprime, multiplication \(M_f:E\to E\) is invertible, and

\[
\eta M_f=M_{f(\sum_i s_i)}\eta.
\]

This supplies the exact arithmetic map behind \(B\). Multiplication by \(f_l\) changes the source representative in a known way; it is not silently identified with the identity. {[}R09 §2, ``The exact two-reference multiplication,'' and §3, ``Joint full jets and the mixed Schur complement.''{]}

Comparison estimates between these norms and the arithmetic source are proved on \emph{all polynomial directions} before taking restrictions or quotients. If \(aH_1\preceq H_2\preceq AH_1\) on one coefficient space, then minimizing each inequality over the same affine fibre proves the corresponding quotient inequality. A bound on selected moments or selected monic test vectors alone would not prove that whole-form comparison. The current arithmetic comparison is AMT, and later source-order refinements retain those same fibres and maps. {[}R14 §8 pp.52ff.; §§192,197,204--211 pp.571ff.,602ff.,648--659.{]}

\section{From a source norm to four determinant returns}\label{from-a-source-norm-to-four-determinant-returns}

For \(N\ge q-1\), the remainder map \(J_N:\mathcal P_N\twoheadrightarrow E\) has kernel \(\chi\mathcal P_{N-q}\). If \(H_N\) is the source Gram in a fixed coefficient frame, its attained quotient Gram is

\[
G_N=(J_NH_N^{-1}J_N^*)^{-1}.
\tag{C}
\]

Indeed \(S_N^{\mathrm{rem}}=H_N^{-1}J_N^*G_N\) is a right inverse of \(J_N\). For \(z\in\ker J_N\), \(z^*H_NS_N^{\mathrm{rem}}=0\). Every lift of a value is the displayed minimum lift plus such a \(z\), so its squared norm is the minimum squared norm plus \(z^*H_Nz\). This proves (C), not merely its positivity. In a monic orthogonal-polynomial frame, with \(b_n=[p_n]_\chi\) and full squared norm \(\omega_n\), the inverse is

\[
G_N^{-1}=\sum_{n=0}^{N}\frac{b_nb_n^*}{\omega_n}.
\]

The first \(q\) remainder columns form a triangular invertible matrix; hence the inverse exists. {[}R14 CAI2, §198 p.622.{]}

Now apply the \emph{actual} observation \(\Lambda:E\twoheadrightarrow B\) and a fixed inclusion \(I:K\hookrightarrow E\). Its kernel metric and attained image metric are

\[
H_{K,N}=I^*G_NI,\qquad
Q_{B,N}=(\Lambda G_N^{-1}\Lambda^*)^{-1},
\] \[
S_N=G_N^{-1}\Lambda^*Q_{B,N},\quad
\Lambda S_N=I_B,\quad I^*G_NS_N=0.
\]

Choose one fixed section \(S_*\) of \(\Lambda\). The frames \([I,S_*]\) and \([I,S_N]\) differ by a triangular matrix with determinant one, since \(S_N-S_*\) has image in \(K\). Orthogonality therefore proves

\[
\det H_{K,N}\det Q_{B,N}
=|\det[I,S_*]|^2\det G_N.
\tag{D}
\]

The frame factor is present; it is not set equal to one. Use the original four endpoints

\[
N_0=q-1,\quad N_1=q,\quad N_2=2q-1,\quad N_3=2q,
\qquad(\varepsilon_0,\varepsilon_1,\varepsilon_2,\varepsilon_3)
=(1,1,-1,-1).
\]

Define the total, kernel and observed-image returns by

\[
\mathcal B_k=\sum_{i=0}^3\varepsilon_i\log\det G_{N_i},\quad
F_K=\sum_i\varepsilon_i\log\det H_{K,N_i},\quad
F_B=\sum_i\varepsilon_i\log\det Q_{B,N_i}.
\]

Taking this exact signed sum in (D) cancels the fixed factor because \(\sum_i\varepsilon_i=0\), and proves \(F_K+F_B=\mathcal B_k\). The signs are part of the definition; they are what makes constant frame changes cancel. This same calculation works for the arithmetic and each Gamma source separately. {[}R14 CAI3--5, pp.622--623; AKP6--10, pp.111--112.{]}

Schur complements, minimum sections, determinant lemmas and their Woodbury form are classical finite-dimensional linear algebra. The contribution to this programme is the calculation of the original coefficient maps, frames and constants to which those identities are applied, and the proof that they give the original arithmetic/cohomological quantities. The classical linear-algebra identities are credited independently of this programme-specific application \cite{BoydElGhaouiFeronBalakrishnan1994,Hager1989Inverse,ShermanMorrison1950Inverse,Woodbury1950Modified}.

\section{The later native metric and its two correlated row energies}\label{the-later-native-metric-and-its-two-correlated-row-energies}

The conductor calculations recast the relevant dual metric in terms of a numerator divided by the actual conductor symbol \(E_A\). Its order \(v=\operatorname{ord}_0E_A\), leading moment \(\mu_v\ne0\) and scalar \(\alpha_A=v!/\mu_v\) remain in every division matrix. The metric is the squared norm of that numerator after minimizing over all allowed lower evaluations. Thus its finite raw data consist of a lower-evaluation matrix \(V\) and a numerator matrix \(T\), both in the original orthonormal Gamma frame.

At one cutoff set \(A=V^*V\), \(K_0=A^{-1}V^*T\). When the next rows \(v_+,t\) are appended, the residual row and its conditioning denominator are

\[
r=t-v_+K_0,\qquad \lambda=1+v_+A^{-1}v_+^*,\qquad
h=\lambda^{-1/2}r.
\]

If the old attained numerator Gram is \(D\), its new value is exactly

\[
D^+=D+h^*h.
\tag{E}
\]

To check \(E\), subtract the old attained minimum. A change \(\eta\) in the old lower coefficient costs

\[
\eta^*A\eta+\|rx-v_+\eta\|^2.
\]

The normal equations give \(\eta=(A+v_+^*v_+)^{-1}v_+^*rx\). Substitution gives \(x^*r^*r x/(1+v_+A^{-1}v_+^*)\). This proves the update with the complete lower conditioning, rather than minimizing the new row independently of the old rows. {[}R14 HJR14--15, §73 pp.304--309.{]}

Let \(B\) now denote the actual onto coefficient map in the native quotient and \(J\) its kernel inclusion. These local symbols belong to HJR and should not be confused with the earlier period-image space \(B\). Set

\[
D_J=J^*DJ,\qquad D_B=(BD^{-1}B^*)^{-1},\qquad
L=D^{-1}B^*D_B.
\]

The same attained-section proof gives the exact orthogonal inverse decomposition

\[
D^{-1}=J D_J^{-1}J^*+L D_B^{-1}L^*.
\]

Consequently define two nonnegative \emph{actual} row energies

\[
x=hJD_J^{-1}J^*h^*,\qquad
y=hLD_B^{-1}L^*h^*.
\]

The determinant lemma and one further quotient minimum give

\[
\log\frac{\det D^+}{\det D}=\log(1+x+y),
\] \[
\log\frac{\det D_J^+}{\det D_J}=\log(1+x),\qquad
\log\frac{\det D_B^+}{\det D_B}
=\log\left(1+\frac{y}{1+x}\right).
\tag{F}
\]

For the last identity, a quotient lift of a value \(a\) is \(La+Jz\) in the stated original coefficient spaces. Its old squared norm is \(a^*D_Ba+z^*D_Jz\), and its additional squared norm is \(|hLa+hJz|^2\). Put \(t=D_J^{1/2}z\) and \(w=hJD_J^{-1/2}\), so \(ww^*=x\). The same one-row normal equation as in (E) gives the minimum additional form \(|hLa|^2/(1+x)\). Its determinant relative to \(D_B\) is \(1+y/(1+x)\), proving the last identity with that denominator. Thus

\[
\log(1+x+y)=\log(1+x)+\log\left(1+\frac{y}{1+x}\right)
\]

is the determinant splitting of this precise metric, not a free numerical partition of a known total. {[}R14 HJR16--18, pp.304--309.{]}

Summing the rows from \(M=q-1\) through (2q-1) with weight one at those two ends and weight two at every interior row gives

\[
S_K=\sum_Mw_M\log\left(1+\frac{y_M}{1+x_M}\right),\qquad
S_R=\sum_Mw_M\log(1+x_M),\qquad
T=S_K+S_R.
\tag{G}
\]

The original dual/coefficient maps connect \(S_K\) to the conductor-transported kernel return; the finite discrepancy from \(F_K\) is explicitly controlled. The complete two-window row total satisfies, on the original simple-quartet fixed-period domain,

\[
\frac{T}{kq}\longrightarrow16C_\partial,
\qquad21.66851180520<16C_\partial<21.66851180521.
\]

The coefficient is the evaluated logarithmic-potential expression

\[
C_\partial=4\log(4/\pi)-10-2F(2,\pi)+4\mathcal L(2,\pi),
\]

where \(F\) and \(\mathcal L\) are the energy and logarithmic moment of the proved equilibrium for \(V_{\alpha,\beta}(x)=\beta\sqrt x-\alpha\log x\). The exact equilibrium and outward certificate are part of the source. For orientation to the classical logarithmic-potential setting, see \cite{SaffTotik1997}; the specific potential, original arithmetic comparison and receiving constants are calculated in R14. This number is the total coefficient. It does not select the separate leading value of either summand in (G). {[}R14 HJR19--25 and subsequent receiving refinements; ECL15--18, §59 pp.243--244; current continuation §``Exact working quantities''.{]}

\section{The precise matrix expression that remains}\label{the-precise-matrix-expression-that-remains}

The latest mixed calculation uses the target degree \(j=k+4\), the original source pairs \((1,1)\) and \((k,j)\), and eight fixed coefficient subspaces in the four-copy target. In the notation of the current prompt, put

\[
I=\operatorname{im}(B_{\rm cof}J_0),\quad
B=\operatorname{im}B_{\rm cof},\quad
S=\operatorname{im}[B_{\rm cof},F_e],\quad J=\ker\mathbb B,
\] \[
K_i=I\cap J,\quad W_i=I+J,\quad W=B+J,\quad
K_c=S\cap W,\quad T_e=S+J.
\]

Here \(F_e\), \(e=0,1\), is the original connecting-map column family at the two low cutoffs. The full connecting constructions and exact coefficient maps are in R14 §§151--159 and §§184--193. The resulting inclusions are the two flags

\[
J\subset W_i\subset W\subset T_e,\qquad
K_i\subset I\subset K_c\subset S.
\]

At each original row let \(D_0\) be its preceding positive numerator form, \(\rho_n\) its lower-evaluation-corrected row and \(\zeta_n\) its complete conditioning scalar. With

\[
\mathcal D_0=\operatorname{diag}(D_0,D_0,D_0,D_0),\qquad
\mathcal U_n=\operatorname{diag}(a_n,a_n,a_n,a_n),
\qquad a_n=\frac{\rho_n}{\sqrt{1+\zeta_n}},
\]

the four-by-four compression associated with a fixed independent frame \(X\) for one of the eight subspaces is

\[
K_{X,n}=\mathcal U_nX(X^*\mathcal D_0X)^{-1}X^*\mathcal U_n^*.
\tag{H}
\]

The empty subspace gives the zero matrix. Formula (H) is unchanged by an invertible change of columns in \(X\). If \(Y=\mathcal U_n\mathcal D_0^{-1/2}\), it is \(Y\Pi_XY^*\), where \(\Pi_X\) is the orthogonal projector onto \(\mathcal D_0^{1/2}\operatorname{im}X\). Hence inclusion of subspaces gives positive matrix order. Also

\[
YY^*=r_nI_4,\qquad
r_n=\frac{\rho_nD_0^{-1}\rho_n^*}{1+\zeta_n},
\qquad0\preceq K_{X,n}\preceq r_nI_4.
\]

The signed threshold profile \(\mathcal P_{k,e}\) takes the four eigenvalues of each matrix, including zeros, applies \(\log\max(1,\lambda)\), and sums over the original rows with their endpoint/interior weights. The positive subspaces are \((T_e,K_c,W_i,K_i)\); the negative ones are \((W,S,J,I)\). The accepted theorem gives

\[
|\mathcal L_e-\mathcal P_{k,e}|
\le\sum_n\vartheta_n\varepsilon(r_n)\le16q_j\log2,
\] \[
\varepsilon(r)=
\begin{cases}
4\log(1+r),&0\le r\le1,\\
4\log\bigl(4r/(1+r)\bigr),&r\ge1.
\end{cases}
\tag{I}
\]

This error is \(o(kq_k)\). It has already been proved for the unchanged original profile. Its proof uses the two nested flags and eigenvalue monotonicity; it assumes no simultaneous diagonalization. {[}R14 NTE1--7, §199 pp.631--632.{]}

The map back to the original allocation is equally important:

\[
4S_{K,j}=\mathcal L_e+V_b+V_{\rm rem}+V_s+V_A.
\tag{J}
\]

The terms are attained quotient returns on the original coefficient spaces: \(V_{\rm rem}\) is the remaining target quotient by \(T_e\); \(V_s\) is the connecting source form on \(S/K_c\); \(V_A\) is the invariant source form on \(I/K_i\) in the target-degree metric; and \(V_b\) is the retained boundary return. Each is computed with the original four endpoints and fixed value frame. These positive returns are correlated with \(\mathcal L_e\), since all use the same original row data. The degree-\(k\) proper-source form has its own centre and comparison map; it is not substituted for \(V_A\). {[}R14 CPQ1--20, §194 pp.587--591; EFP, §§204--211.{]}

Thus the remaining target is to evaluate or rigorously constrain the signed profile in (H)--(I) \emph{together with} the complementary returns in (J), through their actual Gamma recurrence, conductor division and coefficient maps. An arbitrary matrix model, a dimension fraction, or the known total \(T\) does not calculate this expression. Conversely, the fact that the expression is not yet evaluated is not an obstruction theorem for the entire cohomological programme.

\section{Why an allocation result must pass through the whole action}\label{why-an-allocation-result-must-pass-through-the-whole-action}

Return to the complete arithmetic quotient \(E\), its multiplication \(A=M_S\) and its metric \(G_N\). The monic arithmetic recurrence, in the original coordinate \(S=c+iy\), is

\[
Sp_n=p_{n+1}+cp_n-\frac{\omega_n}{\omega_{n-1}}p_{n-1}.
\]

It gives the rank-two identity

\[
AG_N^{-1}+G_N^{-1}A^*-kG_N^{-1}
=\frac{b_{N+1}b_N^*+b_Nb_{N+1}^*}{\omega_N}.
\]

Reflection \(P(S)\mapsto P(k-S)\) is an isometry of the original source and quotient. It forces the total adjacent pairing to vanish. The relative Hermitian action therefore has only two possibly nonzero eigenvalues \(\epsilon_N,-\epsilon_N\). In an orthogonal kernel/image split, the individual mixed pairings need not vanish: they cancel in their sum, and the mixed exterior norm retains the term \(2|z_K|^2\). {[}R14 CAI6--9, p.623.{]}

The exact same-class exterior argument gives

\[
\frac{\delta q(k+1)}2\le\epsilon_N\qquad(N\ge q-1).
\tag{K}
\]

This is a lower benchmark from the hypothetical quartet itself. It is not an independent contradiction. To contradict that quartet by this route, a proved upper control would have to conflict with the lower benchmark on the same original object and regime.

Put \(d_n=\det G_n/\det G_{n-1}\), so \(0<d_n<1\) at the relevant rows. The adjacent formulas retain both factors:

\[
\epsilon_{q-1}^2=\frac{\omega_q}{\omega_{q-1}}(d_q^{-1}-1),
\] \[
\epsilon_N^2=\frac{\omega_{N+1}}{\omega_N}
           (1-d_N)(d_{N+1}^{-1}-1)\quad(N\ge q).
\]

After taking the two original windows, with the original weights, the action is exactly

\[
J_k^{\rm action}
=\mathcal B_k^{\rm ar}+\mathcal W_k^{\rm ar}
 -\mathcal P_- -\mathcal P_+,
\tag{L}
\] \[
\mathcal W_k^{\rm ar}
=\log\frac{\omega_{2q-1}\omega_{2q}}{\omega_{q-1}\omega_q}.
\]

The two nonnegative penalties are the complete weighted sums of \(-\log(1-d_n)\) specified in CAI13. Formula (L) follows by multiplying the adjacent identities: monic-norm ratios telescope, reciprocal \(d_n\) ratios give the four source volumes, and the remaining contraction factors give precisely those penalties. The original kernel and reference boundary contributions cancel through \(F_K+F_B=\mathcal B\). The arithmetic volume return \(\mathcal B_k^{\rm ar}\) remains in (L); it is not a coordinate constant. {[}R14 CAI11--15, §198 pp.624--625.{]}

This explains what a future result can claim. A calculation of the intrinsic mixed allocation is meaningful geometry. Its effect on the original action must be obtained by substitution into (L), including its cancellation and every remaining term. It cannot be promoted to an action upper bound by dropping the term that cancels it.

\section{The growing-degree result already proved}\label{the-growing-degree-result-already-proved}

The preceding 692-page addition concerns the complete relation polynomial. In the real coordinate define

\[
Q(y)=i^{-q}\chi(c+iy),\qquad R_k=k\sqrt{\delta^2+\gamma^2}.
\]

Its roots occur in opposite pairs, with their full multiplicities, and lie in the disk of radius \(R_k\). For every polynomial \(P\) of degree at most \(r\), compare

\[
H_0(P)=\int_{\mathbb R}|y|^{2q}|P(y)|^2\,d\sigma(y),\qquad
H_Q(P)=\int_{\mathbb R}|Q(y)P(y)|^2\,d\sigma(y).
\]

The root-sensitive full-form theorem gives explicit two-sided comparison, including the complete inner region on both half-lines. In its exact notation

\[
\alpha=\pi/2,\quad s_q=2q+1,\quad
Y_{q,r}=\frac{s_q^2}{32\alpha(r+s_q)},\quad
\rho_{q,r}=R_k/Y_{q,r},
\] \[
\eta_{q,r}=\frac{c_+}{c_-}
 \left(1+\frac{4r+4q+2}{\alpha}\right)s_q(3/8)^{s_q},
\quad c_-=\sqrt2e^{-7/6},\quad c_+=\sqrt{2\sqrt5}e^{2/3}.
\]

Under its actual finite guards \(2R_k\le Y_{q,r}\), \(\eta_{q,r}<1\), it proves

\[
(1-\rho_{q,r}^2)^q(1-\eta_{q,r})H_0(P)
\le H_Q(P)
\le\bigl((1+\rho_{q,r}^2)^q+\eta_{q,r}\bigr)H_0(P).
\tag{M}
\]

These are proved conditions with a nonempty verified arithmetic asymptotic range, not missing hypotheses inserted for the desired conclusion. The proof controls all directions through classical Laguerre moment/evaluation bounds \cite{Szego1975,DLMF18Current}, the paired-root estimate and a direct inner-region calculation. It then takes the original monic minima and the original affine quotient minima. The actual relation sequence remains

\[
0\longrightarrow\mathcal P_r\xrightarrow{\times\chi}
\mathcal P_{q+r}\xrightarrow{J_{q+r}}E\longrightarrow0.
\]

The map \(P(S)\mapsto i^{-r}P(c+iy)\) on monic degree-\(r\) polynomials retains its phase and inverse; the entire unit in the arithmetic cohomology map is unchanged. {[}R14 RGC1--8, RGC18--24, §217 pp.685--690; RGI1--3, §215 p.682.{]}

A separately proved all-degree monic reference bound has signal

\[
F_{q,r}=\frac{q^2}{4(q+r)\log(8(q+r)/q)}.
\]

Combining that bound with (M) and the full arithmetic source comparison yields canonical allowance lower bounds for \(r_k=\lfloor qk^\theta\rfloor\), for each fixed \(0<\theta<2/3\) at \(m=1\), and \(0<\theta<4/3\) at each fixed \(m\ge2\). The further logarithmic refinement sets

\[
(d,A_m^{\deg})=
\begin{cases}(2,1),&m=1,\\(3,m-1),&m\ge2,\end{cases}
\quad\theta_*=2(d-1)/3,
\] \[
r_k=\left\lfloor\frac{qk^{\theta_*}}{(\log k)^\zeta}\right\rfloor,
\qquad\zeta>1/3,
\]

and proves

\[
\log\min_{q-1\le N\le q+r_k-1}\epsilon_N
\ge(1-o(1))\frac{A_m^{\deg}}{8\theta_*}
 k^{(d+2)/3}(\log k)^{\zeta-1}.
\tag{N}
\]

The proof retains the floor, the original \(q=e_k(k+1)^2\), the full root radius, and the exact adjacent factors. The new endpoints are \(q-1,q,q+r_k-1,q+r_k\); the old high endpoints cannot be reused after changing \(r\). The root-error-to-signal calculation is \(O((\log k)^{1-3\zeta})+o(1)\), which explains the strict condition \(\zeta>1/3\). That earlier estimate established neither \(\zeta=1/3\) nor the uncorrected pure-power endpoint; the successor in Section~\ref{sec:moving-cutoff} proves a larger range. {[}R14 RGC16--32, pp.688--692; RGI11--14, §213 pp.673--675; RGI8--10, §198 pp.629--630.{]}

These canonical lower bounds eventually exceed the benchmark (K) by a ratio tending to infinity in their stated ranges. They locate the behavior of this specified action construction. They do not evaluate the independent mixed profile in (H)--(J), and they do not close the proposed tau/Deligne programme. A further analytic extension must prove its error for these same forms and compute its effect through (L).

\section{The accepted moving-cutoff successor}
\label{sec:moving-cutoff}

The accepted 726-page edition extends the previous comparison to degrees
whose ratio to the relation degree grows. This extension is useful because
the original question concerns increasingly large spaces of representatives.
The quotient, the complete zero multiplicities, and the arithmetic norm all
stay fixed as defined above. The degree of a representative is allowed to
increase, and the estimates must control that increase uniformly.
Write \textbf{R726} for this edition; its complete new proofs are
GRF1--27 (pp.~698--705), GRE1--43 (pp.~706--712), GEL1--13
(pp.~713--717), and GRA1--36 (pp.~718--726).

The analytic comparison uses the original Gamma measure and classical
Laguerre and elliptic-integral formulas. Their human sources are
\citet{Szego1975}, the authors of DLMF Chapter~18
\citep{DLMF18Current}, and B.~C.~Carlson's DLMF Chapter~19
\citep{DLMF19Current}. The parameter-dependent estimates are proved in
GRF, GRE and GEL; citing those classical formulas does not attribute
the programme's arithmetic transfer to their authors.

\subsection{The scalar ratio measures a constraint on the original source}
For an integer \(j\ge0\), let \(\omega_d^{\rm ar}\) be the least
squared arithmetic norm of a monic degree-\(d\) polynomial, and put
\[
 \nu_j^{\rm ar}=
 \min_{P\text{ monic},\,\deg P=j}\|\chi P\|_{\rm ar}^2,
 \qquad T_j=\nu_j^{\rm ar}/\omega_{q+j}^{\rm ar},
 \qquad \beta_d=\omega_d^{\rm ar}/\omega_{d-1}^{\rm ar}.
\]
Both minima are on their displayed affine spaces with the original
arithmetic norm. Every \(\chi P\) in the first minimum is monic of
degree \(q+j\), but its required divisibility restricts the available
representatives. The ratio records the cost of that restriction.
The quotient determinant calculation gives
\[
 d_{q+j}=T_j^{-1},\qquad
 \epsilon_{q+j-1}^2=\beta_{q+j}(T_j-1)
 \begin{cases}1,&j=0,\\1-T_{j-1}^{-1},&j\ge1.
 \end{cases}
\tag{O}
\]
To see the first identity, use the source frame
\((\chi,\chi S,\ldots,\chi S^j;1,S,\ldots,S^{q-1})\).
Its determinant has modulus one after arranging its monic columns by
degree. Orthogonal elimination therefore factors the source Gram
determinant into the relation Gram determinant and the quotient
determinant. Taking consecutive ratios yields
\(d_{q+j}=\omega_{q+j}^{\rm ar}/\nu_j^{\rm ar}\).
Substitution into the adjacent formulas of the preceding section proves
the second identity, including the first-row exception. The complete
source calculation is R726 GRA9--17, pp.~719--721.

\subsection{The larger degree range and its exact scale}
Retain \(k=4\ell+1\), \(q=[1+k(m-1)](k+1)^2\), and the fixed
complete quartet. Let \(r_k\) be an integer sequence and set
\[
 t_k=r_k/q\longrightarrow\infty,
 \qquad t_k\log t_k=o(q/k),
 \qquad H_k=\frac{q}{t_k\log t_k}.
\]
The accepted result, including the whole earlier band rather than just
its last row, is
\[
 \min_{0\le j\le r_k}\log T_j\sim H_k,
 \qquad
 \log\min_{q-1\le N\le q+r_k-1}\epsilon_N\sim H_k/2.
\tag{P}
\]
Here \(H_k/k\to\infty\). The moving comparison first bounds the
original full forms on every coefficient vector. Taking their minima
on the same monic and quotient fibres then proves the scalar transfer.
Its moving root error and its arithmetic error are controlled throughout
the specified range. Small degrees, compact positive degree ratios, and
growing ratios are joined by the complete uniform argument in R726
GRB5--8, pp.~678--679; the scalar providers are GRA18--25.
Thus (P) is a theorem of the stated source edition, with an explicit
range, not an inference from a numerical fit.

For a concrete consequence, retain the earlier pair
\((d,A_m^{\deg})=(2,1)\) for \(m=1\) and \((3,m-1)\) for fixed
\(m\ge2\). For
\[
 r_k=\left\lfloor\frac{qk^\theta}{(\log k)^\zeta}\right\rfloor,
\]
one obtains
\[
 \log\min_{q-1\le N\le q+r_k-1}\epsilon_N
 \sim\frac{A_m^{\deg}}{2\theta}
      k^{d-\theta}(\log k)^{\zeta-1}
\tag{Q}
\]
when \(0<\theta<d-1\), for any fixed real \(\zeta\), or when
\(\theta=d-1\) and \(\zeta>1\).
Indeed the floor gives
\(t_k=k^\theta(\log k)^{-\zeta}+O(1/q)\), whence
\[
 \frac{t_k\log t_k}{q/k}
 \sim\frac{\theta}{A_m^{\deg}}
       k^{\theta-d+1}(\log k)^{1-\zeta}.
\]
This tends to zero in exactly the displayed cases, and substitution
into (P) proves (Q). The earlier pure-power ranges \(\theta<2/3\)
and \(\theta<4/3\) have therefore been enlarged to \(\theta<1\)
and \(\theta<2\), respectively, with the stated logarithmic
approach to the new endpoints.

\subsection{A separately proved enclosure at the critical coefficient}
The equality case \(\theta=d-1,\zeta=1\) needs its actual
arithmetic error. Define from the original density
\[
 \vartheta_h=\int_{-1}^1w_h(y)\,dy,\qquad
 M_{\pi/2}=\int_{\mathbb R}e^{\pi y/2}w_h(y)\,dy,
 \qquad\mathfrak s_h=\frac\pi2+
       \log\frac{M_{\pi/2}}{\vartheta_h}.
\]
The proved tilted tail bound makes \(M_{\pi/2}\) finite.
Evenness gives
\[
M_{\pi/2}=\int\cosh(\pi y/2)w_h(y)\,dy
\ge\int w_h\ge\vartheta_h>0,
\]
so \(\mathfrak s_h\ge\pi/2\).
For a fixed coefficient satisfying
\[
 0<c_{\deg}<\frac{A_m^{\deg}}{2(d-1)\mathfrak s_h},\quad
 r_k=\left\lfloor\frac{c_{\deg}qk^{d-1}}{\log k}\right\rfloor,
 \quad B=\frac{A_m^{\deg}}{c_{\deg}(d-1)},
\]
R726 GRB11--13, pp.~679--680, proves
\[
 \frac B2-\mathfrak s_h
 \le\liminf\frac1k\log\min_{q-1\le N\le q+r_k-1}\epsilon_N
 \le\limsup\frac1k\log\min_{q-1\le N\le q+r_k-1}\epsilon_N
 \le\frac B2+\mathfrak s_h.
\tag{R}
\]
All limits use the original subsequence \(k=4\ell+1\).
The lower endpoint is positive. This is an enclosure; the earlier
asymptotic equivalent (P) is not asserted at this boundary. The coefficient
\(c_{\deg}\) is independent of \(k\) and is distinct from the
unchanged polynomial centre \(c=k/2\).

\subsection{The complete shifted action keeps every endpoint}
For an integer shift \(s\ge1\), put \(a=q+s-1\), \(b=a+q\),
and \(w_0=w_q=1\), \(w_i=2\) for \(1\le i<q\). Then
\[
 J_{k,s}^{\rm action}=2\sum_{i=0}^q w_i\log\epsilon_{a+i}
 =\log\frac{V_aV_{a+1}}{V_bV_{b+1}}
  +\log\frac{\omega_b\omega_{b+1}}{\omega_a\omega_{a+1}}
  -\mathcal P_{k,s},
\tag{S}
\]
where \(V_N=\det G_N\), \(\omega_N=\omega_N^{\rm ar}\), and
\[
 \mathcal P_{k,s}=
 \sum_{j=0}^{q+1}c_j[-\log(1-T_{s-1+j}^{-1})],\qquad
 (c_0,\ldots,c_{q+1})=(1,3,\underbrace{4,\ldots,4}_{q-2},3,1).
\]
For any positive sequence \(z_N\), collecting coefficients gives
\[
 \sum_{i=0}^q w_i(\log z_{a+i+1}-\log z_{a+i})
 =\log z_b+\log z_{b+1}-\log z_a-\log z_{a+1}.
\]
Apply this to \(V\) and \(\omega\) in (O). The two adjacent
contractions give \(c_j=w_j+w_{j-1}\), with absent weights zero,
and therefore total penalty weight \(4q\). This proves (S).
R726 GRA27--36 and GRB9--10 supply all finite error bounds and the
transfer through the original kernel inclusion and quotient map.

For \(t=s/q\to\infty\) with \(t\log t=o(q/k)\), those bounds give
\[
 J_{k,s}^{\rm action}\sim\frac{2q^2}{t\log t},\qquad
 \log\mathcal P_{k,s}\sim-\frac{q}{t\log t}.
\]
At the critical choice \(s=r_k\) in (R), the separate enclosure is
\[
 2(B-\mathfrak s_h)\le
 \liminf\frac{J_{k,s}^{\rm action}}{qk}\le
 \limsup\frac{J_{k,s}^{\rm action}}{qk}\le2(B+\mathfrak s_h).
\]
These advances concern the complete correlated action. The intrinsic
signed eight-matrix allocation and each CPQ complementary return still
require their own calculation on the maps already defined above.


\section{A usable reading route and credit map}\label{a-usable-reading-route-and-credit-map}

The following route leads from the original objects to their exact maps, the current estimates and the remaining calculation:

{
\begin{longtable}[]{@{}
  >{\raggedright\arraybackslash}p{\dimexpr.5\linewidth-\tabcolsep\relax}
  >{\raggedright\arraybackslash}p{\dimexpr.5\linewidth-\tabcolsep\relax}@{}}
\toprule\noalign{}
\begin{minipage}[b]{\linewidth}\raggedright
Question
\end{minipage} & \begin{minipage}[b]{\linewidth}\raggedright
Exact proof location in the fixed edition
\end{minipage} \\
\midrule\noalign{}
\endhead
\bottomrule\noalign{}
\endlastfoot
What are the quartet, multiplicities and original tensor map? & R09 §1; R14 OCS1--3, §1 pp.20--21; CAI1 and CAI10, §198 pp.622,624. \\
What is the actual period/cyclic observation? & R14 OCS2--21, §1 pp.20--25; OQX3 §3 pp.32--33; PQO/RPC §§4--5 pp.34--42. \\
How is a representative change proved to preserve the cohomology class? & R14 AKP1--5, §26 pp.110--111; full analytic and categorical foundations in Appendix~\ref{app:foundational-maps}. \\
Where are all full-unit and confluent-jet factors retained? & R09 §2; R14 OCS12--18, §1 pp.23--25; RGI1--3 §215 p.682. \\
How are arithmetic and Gamma metrics related? & R09 §§2--3; R14 AMT §8 pp.52ff.; NSO/EOR §197 pp.602--620; EFP §§204--211 pp.648--659. \\
Where are the actual minimum metrics, frames and four signs proved? & R14 AKP6--10 §26 pp.111--112; CAI2--5 §198 pp.622--623; HJR14--22 §73 pp.304--309. \\
What total is known, and why does it not give the two allocations? & R14 ECL15--18 §59 pp.243--244; HJR16--25 §73 pp.304--309; current prompt, ``Exact working quantities.'' \\
Where is the mixed profile connected to the original complex? & R14 §§151--159 pp.466--487; CPR/BTK/RPI/RRC §§200--203 pp.633--647; CPQ §194 pp.587--591. \\
What exactly remains to estimate? & R14 SFB §195 pp.592--598, NTE1--7 §199 pp.631--632, EFP §§204--211; equations (H)--(J) here. \\
Why can an allocation bound cancel in the action? & R14 CAI11--15 §198 pp.624--625; changed-window RGI8--10 pp.629--630. \\
Where are the preceding growing-degree estimates? & R14 RGC1--32 §217 pp.685--692 and RGI1--14. \\
What is the accepted moving-cutoff advance? & R726 GRF/GRE/GEL/GRA pp.698--726 and GRB1--13 pp.631--632,674--675,678--680; Section~\ref{sec:moving-cutoff} here. \\
\end{longtable}
}

\appendix
% Additive human-facing appendix. No private paths or account locators.
% Requires amsmath, amssymb; citation keys are documented in the adjacent key map.
\section{Foundational maps: support, the theta complex, and full arithmetic jets}
\label{app:foundational-maps}

This appendix supplies the elementary constructions underlying the
mathematical route. The order matters: first specify what the scalar and
support labels mean, then construct the cochain complex, then connect its
quotient to the actual completed zeta function and its full local jets.
The distinction between an established human construction and a calculation
made in this programme is maintained throughout. In particular, the theta
quotient belongs to the spectral framework developed by Ralf Meyer,
following Alain Connes; the mathematical route gives its exact comparison
with Meyer's quotient, including the topologies and the factor of two
\cite[\S\S2--4]{Meyer2005}.
The language of blueprints used for the structural base is due to
Oliver Lorscheid \cite{Lorscheid2012Blueprints}.

\subsection{The scalar object and its structural base}
\label{app:scalar-and-base}

The project origin of this construction is documented in
\cite{GlobalizationNote2026,SplitZeroSecondary2026}, work developed under
KokunoYumeto's direction with AI assistance. The deposited creator fields
are retained in the bibliography. Semiring background is credited to
Jonathan S.~Golan \cite{AuditGolan1999}; the monoid-scheme and blueprint
frameworks are credited separately to Anton Deitmar and Oliver Lorscheid
\cite{Deitmar2005,Lorscheid2012Blueprints}. No historical priority claim
for the external-zero construction is made.

Let \(R\) be a nonzero commutative ring with identity. Put
\[
D_R=G(R)=\{\tau\}\sqcup\{r^\bullet:r\in R\},
\qquad e_R=0_R^\bullet.
\]
The two displayed symbols \(D_R\) and \(G(R)\) denote the same object.
Its operations are, for \(r,s\in R\),
\[
r^\bullet+s^\bullet=(r+s)^\bullet,\qquad
r^\bullet s^\bullet=(rs)^\bullet,\qquad
\tau+a=a+\tau=a,\qquad \tau a=a\tau=\tau.
\]
The additive identity is \(\tau\), and the multiplicative identity is
\(1_R^\bullet\). Associativity and commutativity reduce to the ring laws
when all arguments are supported; if an argument of a product is
\(\tau\), that product is \(\tau\). For distributivity, when the
multiplier is \(\tau\) both sides are \(\tau\). When the multiplier is
supported and both summands are supported, distributivity is the ring
identity. If a summand is \(\tau\), the required identity is the additive
identity law. This checks all semiring axioms, including every case
involving absence.

The represented zero and the absent scalar obey
\[
e_R^2=e_R,\qquad e_R+e_R=e_R,\qquad
1_R^\bullet+e_R=1_R^\bullet,\qquad
e_Rr^\bullet=e_R,\qquad e_R\ne\tau.
\]
Let \(\mathbb B=\{0,1\}\) have addition \(1+1=1\) and its usual
multiplication. The scalar and support observations are
\[
\begin{array}{lll}
p_R:D_R\longrightarrow R,&p_R(\tau)=0_R,&p_R(r^\bullet)=r,\\
b_R:D_R\longrightarrow\mathbb B,&b_R(\tau)=0,&b_R(r^\bullet)=1.
\end{array}
\]
Both preserve addition and multiplication by the preceding case
distinction. The map \((p_R,b_R):D_R\to R\times\mathbb B\) is injective:
the second coordinate determines presence, and the first coordinate
then determines the present scalar. In particular, the scalar observation
alone forgets information that the pair retains.

The structural pointed monoid is
\[
\mathbf F_{1,\tau}=\{\tau_0,1\},
\]
with \(1\) an identity and \(\tau_0\) absorbing. Its preaddition is the
relation that the empty sum represents \(\tau_0\), together with the
relations required by the preaddition axioms. No equation
\(1+1=1\) is imposed. Let \(B_R\) be the blueprint associated with
\(D_R\): its pointed multiplicative monoid is the multiplicative monoid of \(D_R\),
and its additive relations are the equalities of finite sums in \(D_R\).
Thus the empty sum represents \(\tau\), while
\(r^\bullet+(-r)^\bullet\) represents \(e_R\).
The map
\[
\jmath_R:\mathbf F_{1,\tau}\longrightarrow B_R,\qquad
\jmath_R(\tau_0)=\tau,\qquad \jmath_R(1)=1_R^\bullet
\]
preserves multiplication, the distinguished elements, and every
generating preaddition relation. It is therefore a blueprint morphism.
Every proper prime ideal contains \(\tau\) and excludes
\(1_R^\bullet\); its inverse image under \(\jmath_R\) is exactly
\(\{\tau_0\}\). This gives the structural map to the one-point
spectrum of \(\mathbf F_{1,\tau}\).

There is also a separate localization
\[
D_R[e_R^{-1}]\cong\mathbb B.
\]
Indeed, invertibility and idempotence force \(e_R=1\) after localization.
The identity \(e_Rr^\bullet=e_R\) then identifies every supported scalar
with \(1\). The homomorphism \(b_R\) sends \(e_R\) to \(1\), so factors
through this localization. Conversely the homomorphism from
\(\mathbb B\) sending \(0\) to the localized \(\tau\) and \(1\) to its
unit is well-defined because the localized unit is additively
idempotent. The two maps are inverse on their displayed generators.
This calculation relates the support observation to the scalar object;
it does not replace the structural morphism \(\jmath_R\).

For the original unital inclusion \(j:\mathbb Z\hookrightarrow\mathbb C\),
the split lift sends \(\tau\) to \(\tau\) and \(n^\bullet\) to
\(j(n)^\bullet\). Direct evaluation on these two cases proves
\[
\begin{array}{ccc}
D_{\mathbb Z}&\xrightarrow{D_j}&D_{\mathbb C}\\
p_{\mathbb Z}\downarrow&&\downarrow p_{\mathbb C}\\
\mathbb Z&\xrightarrow{j}&\mathbb C.
\end{array}
\]
The target of the left scalar observation remains the infinite ring
\(\mathbb Z\).

\subsection{Reconstructing all support fibres}
\label{app:support-equivalence}

The project formulation is given in
\cite[Definition~3.5 and Theorems~3.6--3.7]{SplitZeroSecondary2026}.
Its transport-to-the-join addition is the Płonka-sum operation on the
reduct with addition and supported scalar operations
\citep{TauPlonka1967,TauBonzio2018}; the global zero and its scalar action
are specified explicitly below.
The underlying semilattice decomposition belongs to the tradition of
Clifford's semilattices of groups \cite{TauClifford1941}; inverse-semimodule
characterizations are also treated by Sen, Bhuniya and Maity
\cite{TauSenBhuniyaMaity2021}. The source guide distinguishes that
authenticated antecedent from a complete comparison of its hypotheses.
Here all maps for the stated scalar object are proved explicitly.

A diagram consists of a join-semilattice \(L\) with bottom \(\bot\),
an \(R\)-module \(V_l\) for every \(l\in L\), and linear maps
\(\rho_{lk}:V_l\to V_k\) for \(l\le k\), satisfying
\(\rho_{ll}=1\) and \(\rho_{kj}\rho_{lk}=\rho_{lj}\).
A morphism \((L,V)\to(L',W)\) is a bottom- and join-preserving map
\(\alpha:L\to L'\), together with linear maps
\[
f_l:V_l\longrightarrow W_{\alpha(l)},\qquad
f_k\rho_{lk}=\rho'_{\alpha(l),\alpha(k)}f_l.
\]
Define its reconstructed carrier and operations by
\[
\mathcal T(L,V)=\coprod_{l\in L}\{l\}\times V_l,
\]
\[
(l,x)+(k,y)=
\bigl(l\vee k,\rho_{l,l\vee k}x+\rho_{k,l\vee k}y\bigr),
\quad
\tau(l,x)=(\bot,0),\quad
r^\bullet(l,x)=(l,rx).
\]
Its global zero is \((\bot,0)\). No condition \(V_\bot=0\) is imposed.
Threefold addition transports both parenthesizations to
\(l\vee k\vee j\), where their amplitudes are
\(\rho_lx+\rho_ky+\rho_jz\). This proves associativity.
Commutativity and the bottom identity follow from the same transport
formulas. Supported scalar laws follow from the module laws in each
fibre. In particular \(e_R(l,x)=(l,0)\). Distributivity over vectors
follows from linearity of the transition maps. Distributivity when a
scalar is absent uses the global bottom zero; all remaining cases
reduce to the ring laws. Hence \(\mathcal T(L,V)\) is a
\(D_R\)-semimodule. The formula
\[
\mathcal T(\alpha,f)(l,x)=(\alpha(l),f_lx)
\]
is linear: the displayed compatibility of the \(f_l\) proves its
addition law after transport to the join, and the scalar laws hold in
the indicated fibre.

Conversely, let \(M\) be a \(D_R\)-semimodule. Define
\[
L_M=e_RM,\qquad l\vee k=l+k,\qquad \bot=0_M,
\]
\[
V_l=\{m\in M:e_Rm=l\},\qquad
\rho_{lk}(m)=m+k\quad\text{when }l+k=k.
\]
The equality \(e_R^2=e_R\) makes every \(l\in L_M\) fixed by \(e_R\).
The equality \(e_R+e_R=e_R\) makes \(l+l=l\). Thus addition in \(L_M\)
is a commutative idempotent monoid, giving the stated join order and
bottom. If \(m,n\in V_l\), then \(e_R(m+n)=l+l=l\), so this fibre is
closed under addition. Moreover
\[
m+l=(1_R^\bullet+e_R)m=m,\qquad
m+(-1_R)^\bullet m=e_Rm=l.
\]
Thus \(V_l\), with zero \(l\), is an abelian group. For a supported
scalar, \(e_R(r^\bullet m)=l\); the inherited ring-module laws,
including \(0_Rm=l\), make it an \(R\)-module.
If \(l\le k\), then \(e_R(m+k)=l+k=k\). The map \(m\mapsto m+k\)
sends zero to zero and is additive because \(k+k=k\). It is
\(R\)-linear because \(r^\bullet k=k\). Its identity and composition
laws follow from \(m+l=m\) and \(k+j=j\).

A linear map \(F:M\to N\) sends \(e_RM\) to \(e_RN\) and induces
\(\alpha(l)=F(l)\). This preserves bottom and joins and restricts to
linear maps of the displayed fibres. Naturality is the equality
\(F(m+k)=F(m)+F(k)\). These constructions are inverse functors up to
the following explicit isomorphisms:
\[
\mathcal T(L_M,V)\longrightarrow M,\quad(l,m)\longmapsto m,
\qquad
M\longrightarrow\mathcal T(L_M,V),\quad m\longmapsto(e_Rm,m),
\]
\[
L\longrightarrow e_R\mathcal T(L,V),\quad l\longmapsto(l,0),
\qquad V_l\longrightarrow\{l\}\times V_l,\quad x\longmapsto(l,x).
\]
The first pair are inverse as functions. Their addition law reduces
to
\[
(m+(l\vee k))+(n+(l\vee k))
=m+n+(l\vee k)=m+n,
\]
since \(e_R(m+n)=l\vee k\); their scalar laws follow directly.
The second pair preserve the join and every module and transition
map by their definitions. All formulas commute with \(F\) or
\((\alpha,f)\), proving naturality and hence the asserted equivalence.

For later use, a split-linear map \(F=\mathcal T(\alpha,f)\) has
two precisely related zero conditions:
\[
F(l,x)=0_N
\iff \alpha(l)=\bot'\ \text{and}\ f_lx=0,\qquad
F(l,x)\in e_RN\iff f_lx=0.
\]
If \(L_0=\alpha^{-1}(\bot')\), their carriers are respectively
\[
\mathcal T\bigl(L_0,(\ker f_l)_{l\in L_0}\bigr)
\ \longrightarrow\
\mathcal T\bigl(L,(\ker f_l)_{l\in L}\bigr),
\]
where the arrow is inclusion. The set \(L_0\) is closed under joins.
The transition maps preserve these kernels because
\(f_k\rho_{lk}=\rho'_{\alpha(l),\alpha(k)}f_l\).
Any linear map satisfying either zero condition therefore factors
uniquely through the corresponding inclusion. The second carrier is
the fibre product \(M\times_N e_RN\). This proves the universal
properties as well as the exact comparison between them.

\subsection{The four-point chart and its derived global sections}
\label{app:chart-resolution}

Let \(P\) be the poset
\[
+<\eta<\sigma,\qquad -<\eta<\sigma,
\]
where \(+\) and \(-\) are incomparable. Give it the upper-set
topology and work in the algebraic coefficient category
\(\mathcal A=\operatorname{Fun}(P,\operatorname{Vect}_{\mathbb C})\)
of covariant diagrams. For a point \(x\), let \(P_x(y)=\mathbb C\)
when \(x\le y\), and zero otherwise, with identity maps on nonzero
arrows. A natural transformation \(P_x\to F\) is determined by the
image of \(1\) at \(x\), and at \(y\ge x\) that image must be
\(r_{xy}v\). Hence
\[
\operatorname{Hom}_{\mathcal A}(P_x,F)\cong F_x.
\]
Evaluation is exact, since diagram kernels and cokernels are
computed at each point. Thus \(P_x\) is projective.

Let \(\underline{\mathbb C}\) be the constant diagram. The sequence
\[
0\longrightarrow P_\eta
\xrightarrow{v\mapsto(v,-v)}P_+\oplus P_-
\xrightarrow{(u,v)\mapsto u+v}\underline{\mathbb C}
\longrightarrow0
\]
is exact. At \(+\) or \(-\), its nonzero part is
\(\mathbb C\xrightarrow{1}\mathbb C\). At \(\eta\) and \(\sigma\),
its nonzero part is
\(\mathbb C\xrightarrow{(1,-1)}\mathbb C^2
\xrightarrow{(1,1)}\mathbb C\); the kernel of addition is the
displayed diagonal with opposite signs. All restriction squares
commute with the component maps, proving exactness as diagrams.

Writing \(r_\pm:F_\pm\to F_\eta\), compatible global sections are
pairs \((u,v)\) with \(r_+u=r_-v\); their values at \(\eta\) and
\(\sigma\) are forced. Thus
\(\Gamma(P,F)=\operatorname{Hom}_{\mathcal A}
(\underline{\mathbb C},F)\), and the projective resolution gives
\[
\mathsf A_\tau(F)=R\Gamma(P,F)\simeq C_F,\qquad
C_F^0=F_+\oplus F_-,\quad C_F^1=F_\eta,\quad
d_F(u,v)=r_+u-r_-v.
\]
The cohomological degrees and sign in this formula are fixed.
With \(r=r_{\eta\sigma}\), the two maps \(C_F\to F_\sigma[0]\)
have degree-zero parts
\[
c_+(u,v)=rr_+u,\qquad c_-(u,v)=rr_-v.
\]
Both are chain maps. Their difference is the chain homotopy
\(h^1=r\), since \(c_+-c_-=h^1d_F\), and all degree-one target
components vanish. They therefore represent a single natural
derived restriction morphism.

\subsection{The algebraic right adjoint, including the signs}
\label{app:chart-adjunction}

For a vector space \(W\), define \(I_x(W)_y=W\) if \(y\le x\)
and zero otherwise. A morphism \(F\to I_x(W)\) is uniquely
determined by \(\ell:F_x\to W\): at \(y\le x\) it is
\(\ell r_{yx}\), and elsewhere it is zero. Hence
\[
\operatorname{Hom}_{\mathcal A}(F,I_x(W))
\cong\operatorname{Hom}_{\mathbb C}(F_x,W).
\]
Every vector space is injective as a \(\mathbb C\)-module: extend
a basis of a subspace to a basis and choose arbitrary images for
the extra basis vectors. Since evaluation is exact, \(I_x(W)\)
is injective in \(\mathcal A\).

Put
\[
K_\tau(W)^{-1}=I_\eta(W),\qquad
K_\tau(W)^0=I_+(W)\oplus I_-(W),\qquad
d_K=(+\mathrm{res},-\mathrm{res}).
\]
It is a bounded complex of injectives. Its Hom complex from a
coefficient sheaf \(F\) is
\[
\left[
\operatorname{Hom}_{\mathbb C}(F_\eta,W)
\xrightarrow{\ell\mapsto(\ell r_+,-\ell r_-)}
\operatorname{Hom}_{\mathbb C}(F_+,W)\oplus
\operatorname{Hom}_{\mathbb C}(F_-,W)
\right]
\]
in degrees \(-1,0\). In the convention
\[
d_{\operatorname{Hom}}\alpha
=d_W\alpha-(-1)^{|\alpha|}\alpha d_F,
\]
the complex \(\operatorname{Hom}^\bullet(C_F,W[0])\) has
exactly these terms and this differential. Consequently there
is a natural chain isomorphism
\[
R\operatorname{Hom}_{\mathcal A}(F,K_\tau(W))
\cong R\operatorname{Hom}_{\mathbb C}(\mathsf A_\tau(F),W).
\]

For a bounded complex \(W^\bullet\), an explicit extension is
\[
K_\tau(W)^n=
I_\eta(W^{n+1})\oplus I_+(W^n)\oplus I_-(W^n),
\]
\[
d(s,u,v)=
\bigl(-d_Ws,\ \mathrm{res}_+s+d_Wu,\
-\mathrm{res}_-s+d_Wv\bigr).
\]
Its square is zero: the first component is \(d_W^2s\);
in each other component the two restriction-of-\(d_Ws\) terms
cancel, leaving \(d_W^2u\) or \(d_W^2v\).
For a sheaf \(F\), the degree-\(n\) Hom comparison sends
\[
(s,u,v)\longmapsto
\bigl(a=(u,v),\ b=(-1)^{n+1}s\bigr)
\in\operatorname{Hom}^n(C_F,W).
\]
Indeed its differential on \(C_F^0\) is
\[
d_Wa-(-1)^n b\,d_F=d_Wa+s\,d_F,
\]
which gives the two restriction terms above, and its differential
on \(C_F^1\) is
\((-1)^{n+1}d_Ws\), equal to the image of \(-d_Ws\) under
the degree-\(n+1\) comparison. This verifies the phase explicitly.
The same identification for bounded coefficient complexes \(F\)
is obtained by finite totalization with this Hom differential;
the component identities just checked are its horizontal
identities, while naturality makes it commute with the
differential of \(F\).

At \(+\), \(-\), \(\eta\), and \(\sigma\), the stalk complexes for
\(W[0]\) are respectively
\[
[W\xrightarrow{1}W],\quad
[W\xrightarrow{-1}W],\quad W[1],\quad0.
\]
Let \(S_\eta W\) be the diagram supported only at \(\eta\).
The map \(S_\eta W[1]\to K_\tau(W)\), identity at \(\eta\)
and zero at other stalks, is natural and a quasi-isomorphism:
the first two displayed stalk complexes are acyclic, and the
remaining stalk comparisons are identities or zero-to-zero.
Thus a zero stalk at \(\sigma\) does not make this right adjoint
the zero object. All assertions here concern the stated
algebraic diagram category. This calculation alone asserts
neither a continuous-dual identification for a Schwartz
quotient nor an involutive Verdier duality on an arithmetic
scheme.

\subsection{The original analytic source and the two-leg complex}
\label{app:theta-source}

The original Fourier convention and source are
\[
\widehat\phi(\xi)=\int_{\mathbb R}\phi(x)e^{-2\pi ix\xi}\,dx,
\qquad
V=\{\phi\in\mathcal S(\mathbb R):
\phi(-x)=\phi(x),\ \phi(0)=\widehat\phi(0)=0\}.
\]
On the positive half-line retain
\[
D=-x\partial_x,\qquad
\mathscr B=\left\{F\in C^\infty(\mathbb R_{>0}):
p_{b,j}(F)=\sup_{x>0}x^b|D^jF(x)|<\infty\
(b\in\mathbb Z,\ j\ge0)\right\},
\]
\[
\Theta\phi(x)=\sum_{n\in\mathbb Z\setminus\{0\}}\phi(nx)
=2\sum_{n\ge1}\phi(nx),\qquad JF(x)=x^{-1}F(1/x).
\]
Poisson summation is the classical input recorded also in
Meyer's construction \cite[\S3, equation (11)]{Meyer2005}.
In this convention it follows by periodizing a Schwartz function:
\(\sum_{n\in\mathbb Z}\phi(y+nx)\) is a smooth \(x\)-periodic
function whose \(k\)-th Fourier coefficient is
\(x^{-1}\widehat\phi(k/x)\), by a change of variables on each
interval of the periodization. The original series and all
its derivatives converge uniformly on a period by Schwartz
decay; the Fourier coefficients decrease faster than every
power, by repeated integration by parts. Its Fourier series
therefore converges absolutely to the periodized function.
Evaluation at \(y=0\) gives
\[
\sum_{n\in\mathbb Z}\phi(nx)
=x^{-1}\sum_{k\in\mathbb Z}\widehat\phi(k/x).
\]
For \(\phi\in V\), removal of both zero terms yields
\[
\Theta\phi(x)=x^{-1}\Theta\widehat\phi(1/x),
\qquad J\Theta=\Theta\mathcal F.
\]
Fourier inversion gives \(\mathcal F^2\phi(x)=\phi(-x)\);
hence \(\mathcal F^2=1\) on even functions, and it exchanges
the two vanishing moments. Thus it preserves \(V\).

Here are the actual estimates ensuring that the theta operator
has the stated target. For \(N>1\), \(x\ge1\), and \(j\ge0\),
\[
|D^j\Theta\phi(x)|
\le2\zeta(N)x^{-N}\sup_{y>0}y^N|D^j\phi(y)|.
\]
Each sum converges normally on compact positive intervals.
The identity \(DJ=J(1-D)\), checked by differentiation, and
the preceding Poisson identity give at \(0<x\le1\) a finite
sum of the same bounds for \(\widehat\phi\) at \(1/x\).
For every desired \(b,j\), choose \(N\) large enough that
\(N+b-1\ge0\) for the latter bounds and \(N-b\ge0\) for the
former. They bound \(p_{b,j}(\Theta\phi)\) by finitely many
Schwartz seminorms of \(\phi\) and \(\widehat\phi\).
The Fourier transform is continuous on the Schwartz space,
as follows by differentiating its integral and integrating
by parts. This proves continuity \(V\to\mathscr B\).

The map \(\Theta\) is injective. For fixed \(x>0\), choosing
a Schwartz exponent \(N>2\) makes
\(\sum_{n,r\ge1}|\phi(nrx)|\) converge: it is bounded by a
constant times \(\sum_{n,r}(nr)^{-N}\).
The divisor identity for the Möbius function, obtained by
expanding \(\prod_{p\mid d}(1-1)\), therefore gives
\[
\frac12\sum_{n\ge1}\mu(n)\Theta\phi(nx)
=\sum_{d\ge1}\left(\sum_{n\mid d}\mu(n)\right)\phi(dx)
=\phi(x).
\]
Evenness and the prescribed value at zero then recover
\(\phi\) on the full real line.

For the diagram
\(\mathcal T=(V,V,\mathscr B,0)\) with restrictions
\(\Theta\) and \(J\Theta\), the chart resolution is the
literal complex
\[
C_\theta=[V\oplus V\xrightarrow{d}\mathscr B],\qquad
d(\phi,\psi)=\Theta(\phi-\widehat\psi),
\]
in degrees \(0,1\). Its image is \(\Theta V\), since
Fourier preserves \(V\) and every \(\Theta\phi\) is
\(d(\phi,0)\). Injectivity of \(\Theta\) consequently proves
\[
H^0(C_\theta)=\{(\widehat\psi,\psi):\psi\in V\},
\qquad H^1(C_\theta)=Q=\mathscr B/\Theta V.
\]
For a singleton mask the differential is respectively
\(\Theta\) or \(-\Theta\mathcal F\). Both are injective
and have image \(\Theta V\), giving zero degree-zero
cohomology and the same degree-one quotient \(Q\).
Coordinate inclusion of that leaf and the identity on
\(\mathscr B\) form a chain map to the two-leg complex;
on \(H^1\) it is the identity of \(Q\). Reconstructing the
full support diagram gives
\[
\widetilde H^1_\tau
=\{\tau\}\sqcup
\coprod_{\varnothing\ne A\subseteq\{+,-\}}\{A\}\times Q.
\]
The skeleton map is
\(\tau\mapsto\tau\), \((A,[F])\mapsto(A,0)\).
Its inverse image of the global absent zero is
\(\{\tau\}\); its inverse image of the whole supported
skeleton is the entire carrier. These assertions follow
from the two zero conditions already proved in
Section~\ref{app:support-equivalence}.

\subsection{The Gaussian, its Mellin transform, and every original constant}
\label{app:gaussian-mellin}

Put
\[
u(x)=e^{-\pi x^2},\qquad
\phi_*(x)=(4\pi^2x^4-6\pi x^2)e^{-\pi x^2}.
\]
The positive integral \(I=\int_{\mathbb R}u(x)\,dx\)
satisfies
\[
I^2=\int_{\mathbb R^2}e^{-\pi(x^2+y^2)}\,dx\,dy
=2\pi\int_0^\infty e^{-\pi r^2}r\,dr=1,
\]
so \(I=1\). Differentiation under the Fourier integral
and integration by parts, justified by Gaussian decay,
give \(\partial_\xi\widehat u=-2\pi\xi\widehat u\).
Its initial value is \(\widehat u(0)=1\), hence
\(\widehat u(\xi)=e^{-\pi\xi^2}\).
The same operations for a Schwartz function give
\[
\mathcal F(xf)=\frac{i}{2\pi}\partial_\xi\widehat f,
\qquad \mathcal F(f')=2\pi i\xi\widehat f,
\qquad \mathcal F D=(1-D)\mathcal F.
\]
Direct differentiation gives
\[
Du=2\pi x^2u,\qquad
D^2u=(4\pi^2x^4-4\pi x^2)u,\qquad
\phi_*=D(D-1)u.
\]
Consequently
\[
\widehat{\phi_*}
=(1-D)(-D)u=D(D-1)u=\phi_*.
\]
This proves that \(\phi_*\) is even Schwartz,
\(\phi_*(0)=0\), and
\(\widehat{\phi_*}(0)=0\), so \(\phi_*\in V\).
The operator \(D\) preserves \(V\): it preserves
Schwartz regularity and evenness, has value zero at zero,
and integration by parts gives
\(\int_{\mathbb R}D\phi=\int_{\mathbb R}\phi=0\)
when \(\phi\in V\). Hence every polynomial in \(D\)
used below remains on the original source.

For \(F\in\mathscr B\), define
\[
\mathcal MF(s)=\int_0^\infty F(x)x^s\,\frac{dx}{x}.
\]
On every compact set of \(s\), choose integers strictly
below and strictly above all occurring real parts.
The two corresponding seminorms bound the integral
at zero and infinity, including every power of
\(\log x\). Differentiation under the integral
therefore proves that \(\mathcal MF\) is entire.
Integration by parts, with vanishing endpoint term,
gives
\[
\mathcal M(DF)(s)=s\,\mathcal MF(s).
\]
For the positive-half-line transform \(M_+\phi_*\),
the substitution \(t=\pi x^2\) gives, initially for
\(\Re s>0\),
\[
\begin{aligned}
M_+\phi_*(s)
&=\pi^{-s/2}
\left(2\Gamma\left(\frac s2+2\right)
-3\Gamma\left(\frac s2+1\right)\right)\\
&=\frac12s(s-1)\pi^{-s/2}\Gamma(s/2).
\end{aligned}
\]
The recurrence used here is obtained by integration
by parts in the defining Gamma integral. For
\(\Re s>1\), absolute convergence permits summation
inside the theta integral, since
\[
\sum_{n\ne0}\int_0^\infty
|\phi_*(nx)|x^{\Re s}\,\frac{dx}{x}
=2\zeta(\Re s)\int_0^\infty
|\phi_*(y)|y^{\Re s}\,\frac{dy}{y}<\infty.
\]
Therefore, with \(f_0=\Theta\phi_*\),
\[
\mathcal Mf_0(s)
=2\zeta(s)M_+\phi_*(s)
=s(s-1)\pi^{-s/2}\Gamma(s/2)\zeta(s)
=g(s)=2\xi(s).
\]
Here \(\xi(s)=\tfrac12s(s-1)\pi^{-s/2}
\Gamma(s/2)\zeta(s)\) is the programme's Riemann
convention. The factor two from the integer sum
and the factor one half in \(M_+\phi_*\) both remain.
The entire function \(\mathcal Mf_0\) provides the
continuation in this identity. Also \(Jf_0=f_0\),
so the substitution \(x\mapsto1/x\) proves
\(g(1-s)=g(s)\). Reality of \(f_0\) gives
\(\overline{g(\bar s)}=g(s)\).

For an arbitrary \(\phi\in V\), the same absolutely
convergent calculation gives
\[
\mathcal M\Theta\phi(s)=2\zeta(s)M_+\phi(s).
\]
Since evenness and \(\phi(0)=0\) imply
\(\phi(x)=O(x^2)\) at zero, \(M_+\phi\) is
holomorphic for \(\Re s>-2\). In particular, on
the strip \(0<\Re s<1\) set
\[
H_\phi(s)=
\frac{2\pi^{s/2}M_+\phi(s)}
{s(s-1)\Gamma(s/2)}.
\]
It is holomorphic there and continuation of the
preceding identity gives
\[
\mathcal M\Theta\phi=gH_\phi,\qquad
H_{\phi_*}=1,\qquad H_{P(D)\phi_*}(s)=P(s).
\]
For the last identity, integration by parts on the
positive half-line gives \(M_+(D\phi)=sM_+\phi\)
on this strip; its endpoint term vanishes.

\subsection{Euler inversion and a source section for all local jets}
\label{app:euler-source-section}

Let \(Z\) be a nonempty finite set of actual zeros of
\(g\) in \(0<\Re s<1\), each retained with its full
order \(m_\rho\). Define
\[
h(X)=\prod_{\rho\in Z}(X-\rho)^{m_\rho},
\qquad d_h=\deg h,\qquad E_h=\mathbb C[X]/(h),
\qquad A=M_X.
\]
The jet map \(j_h\) assigns the Taylor class through
order \(m_\rho-1\) at every \(\rho\).
Here is an explicit justification of its polynomial
coordinate interpretation. If
\(H_\rho(X)=h(X)/(X-\rho)^{m_\rho}\), then a
prescribed germ \(f_\rho\) in the \(\rho\)-block
is inserted, with all other blocks zero, by
\[
\left[
H_\rho(X)\sum_{j=0}^{m_\rho-1}\frac1{j!}
\left.\frac{d^j}{dX^j}
\left(\frac{f_\rho(X)}{H_\rho(X)}\right)\right|_{X=\rho}
(X-\rho)^j
\right]_h.
\]
At \(\rho\), the Taylor truncation has exactly the
prescribed jet. At any other selected point,
\(H_\rho\) vanishes to the full required order.
The formula is unchanged if \(f_\rho\) changes by
a germ of order at least \(m_\rho\).
Summing these insertions realizes every jet packet.
A polynomial has zero packet exactly when it is
divisible by every \((X-\rho)^{m_\rho}\), hence by
their product \(h\). This proves
\[
E_h\cong\bigoplus_{\rho\in Z}
\mathcal O_\rho/((s-\rho)^{m_\rho})
=\bigoplus_{\rho\in Z}\mathcal O_\rho/(g),
\]
with the actual coordinates and orders unchanged.

Put
\[
v_h=g/h,\qquad
\varepsilon=j_h(h/g)\in E_h^\times,\qquad
R_Z(u)=\operatorname{rem}_h(\varepsilon u).
\]
Both quotients in these expressions are taken as
germs at the selected zeros; complete multiplicity
makes them holomorphic units there.
In particular \(j_h(v_h)\varepsilon=1\).
Define \(J_Z=j_h\mathcal M:\mathscr B\to E_h\).
The factorization \(gH_\phi\) proves
\(J_Z\Theta=0\), with all jets, not only values.
Thus \(J_Z\) descends to \(\overline J_Z:Q\to E_h\).

We now prove the inverse used for a section, including
its analytic domain. For \(F\in\mathscr B\) with
\(\mathcal MF(\rho)=0\), put
\[
S_\rho F(x)=x^{-\rho}
\int_x^\infty F(y)y^{\rho-1}\,dy.
\]
All complex powers use the real logarithm on
\(\mathbb R_{>0}\). Differentiating gives
\((D-\rho)S_\rho F=F\). The zero moment rewrites
the upper integral as minus the integral from zero
to \(x\). If \(N>\Re\rho\) and \(N+\Re\rho>0\),
these two formulas give respectively
\[
|S_\rho F(x)|\le
\frac{p_{N,0}(F)}{N-\Re\rho}x^{-N},
\qquad
|S_\rho F(x)|\le
\frac{p_{-N,0}(F)}{N+\Re\rho}x^N.
\]
Use the first for \(x\ge1\), the second for
\(0<x\le1\), and enlarge \(N\) for each requested
weight. The identity \(DS_\rho F=\rho S_\rho F+F\)
expresses every Euler derivative of \(S_\rho F\)
as \(\rho^jS_\rho F\) plus a finite linear
combination of Euler derivatives of \(F\).
This proves \(S_\rho F\in\mathscr B\) and
continuity on the specified closed moment kernel.
The homogeneous solution \(cx^{-\rho}\) belongs
to \(\mathscr B\) only if \(c=0\), by considering
both ends and all integer weights. Thus this
inverse is unique. Taking Mellin transforms gives
\[
\mathcal MS_\rho F(s)=\frac{\mathcal MF(s)}{s-\rho},
\]
with the removable value filled in.

Iterate once for each prescribed vanishing factor.
Division lowers that zero order by one and preserves
all other required zeros. This constructs the
continuous inverse \(S_h\) of \(h(D)\) on
\(\ker J_Z\). Conversely, the Mellin transform of
\(h(D)F\) is \(h(s)\mathcal MF(s)\), so lies in that
kernel. Each \(D-\rho\) is injective on \(\mathscr B\)
by the homogeneous-solution calculation. Hence
\[
h(D)\mathscr B=\ker J_Z,\qquad
h(D):\mathscr B\longrightarrow\ker J_Z
\text{ is a continuous bijection with inverse }S_h.
\]

Define the actual source section
\[
R:E_h\longrightarrow\mathscr B,\qquad
Ru=S_h\Theta\bigl(R_Z(u)(D)\phi_*\bigr).
\]
Its input has Mellin transform \(gR_Z(u)\), hence
belongs to the full required jet kernel. The
inverse calculation proves
\[
\mathcal MRu=v_hR_Z(u),\qquad
J_ZRu=j_h(v_h)\varepsilon u=u.
\]
Thus, for \(q_Q:\mathscr B\to Q\),
\[
\sigma=q_QR:E_h\hookrightarrow Q,\qquad
\overline J_Z\sigma=1,\qquad h(D)\sigma=0.
\]
The last equality follows because
\(h(D)Ru=\Theta(R_Z(u)(D)\phi_*)\), an original
theta boundary. The left inverse proves injectivity.
This section contains the full arithmetic unit
\(\varepsilon\); it is not obtained by setting
that unit to \(1\).

\subsection{The exact scaling defect before quotienting}
\label{app:scaling-defect}

Set \(\ell(u)=[X^{d_h-1}]R_Z(u)\).
The polynomial
\(XR_Z(u)-R_Z(Au)\) vanishes modulo \(h\), has
degree at most \(d_h\), and has leading coefficient
\(\ell(u)\). Since \(h\) is monic, this proves
\[
XR_Z(u)-R_Z(Au)=h(X)\ell(u).
\]
The operators \(D\) and \(\Theta\) commute by
termwise differentiation. Apply \(h(D)\) to
\(DR-RA\), use the preceding polynomial identity,
and compare with \(h(D)f_0\ell\). Injectivity of
\(h(D)\) proves the full source identity
\[
DR-RA=f_0\ell.
\]
For one simple zero this coefficient is
\(\ell(u)=u/g'(\rho)=u/(2\xi'(\rho))\).

For \(a>0\), let \(U_aF(x)=F(x/a)\).
It preserves both original function spaces and
satisfies \(U_a\Theta=\Theta U_a\).
Substitution in the Fourier integral gives
\(\mathcal F U_a=aU_{1/a}\mathcal F\).
Consequently the original chain action is
\[
\mathcal U_a^0(\phi,\psi)
=(U_a\phi,aU_{1/a}\psi),\qquad
\mathcal U_a^1F=U_aF,
\]
and \(d\mathcal U_a^0=\mathcal U_a^1d\).

For \(t\in\mathbb R\), put \(T_t=e^{tA}\) and
\[
c_t(u)=\int_0^t U_{e^{t-r}}\phi_*\,\ell(T_ru)\,dr.
\]
The integral is oriented, so it is negative the
integral from \(t\) to zero when \(t<0\).
For the Schwartz seminorm
\(q_{N,j}(\phi)=\sup_x(1+|x|)^N|\partial_x^j\phi(x)|\),
direct substitution gives
\[
q_{N,j}(U_b\phi)\le
b^{-j}\max(1,b)^Nq_{N,j}(\phi).
\]
Using any fixed norm on the finite-dimensional
\(E_h\), only for this estimate, yields
\[
q_{N,j}(c_tu)\le
|t|e^{(N+j+\|A\|)|t|}
q_{N,j}(\phi_*)\|\ell\|\,\|u\|.
\]
The integrand is continuous in each seminorm.
Its Riemann sums therefore converge in the complete
closed Schwartz subspace \(V\), proving that the
integral is an actual vector of that source.
The analogous identity
\(p_{b,j}(U_aF)=a^bp_{b,j}(F)\) controls the
target space. Differentiation of dilations in each
of these seminorms follows from the integral
remainder for the derivative and the same bounds
applied to \(D^2\phi\) or \(D^2F\).
Thus, in \(\mathscr B\),
\[
\frac{d}{dr}\bigl(U_{e^{t-r}}RT_ru\bigr)
=-U_{e^{t-r}}(DR-RA)T_ru
=-\Theta U_{e^{t-r}}\phi_*\,\ell(T_ru).
\]
Integration proves, including the sign for negative \(t\),
\[
U_{e^t}R-RT_t=\Theta c_t.
\]
Inserting \(U_{e^t}RT_s\) into the defect for
\(t+s\), and then using injectivity of \(\Theta\),
proves the cocycle identity
\[
c_{t+s}=U_{e^t}c_s+c_tT_s,\qquad c_0=0.
\]
In the two-leg complex the two literal homotopies
are
\[
H_t^+(u)=(c_t(u),0),\qquad
H_t^-(u)=(0,-\widehat{c_t(u)}).
\]
Both have differential \(\Theta c_t\).
Their difference
\((c_t,\widehat c_t)\) is the degree-zero cycle
obtained from the Fourier graph; it has not been
discarded. After quotienting,
\[
\overline U_a\sigma=\sigma a^A\qquad(a>0).
\]
This proves equivariance of the section in \(Q\)
and states exactly the retained boundary at the
source level.

\subsection{The cited finite-field theorem and the full-jet comparison}
\label{app:deligne-comparison}

For clarity, the classical theorem being used as
an inspiration has specific geometric hypotheses.
Let \(\mathbb F_q\) be a finite field of characteristic
\(p\), let \(X_0\) be smooth and projective over it,
let \(X=X_0\times_{\mathbb F_q}\overline{\mathbb F}_q\),
and let \(\ell\ne p\) be prime. Deligne's theorem
states that the Frobenius characteristic polynomial
on \(H^i_{\text{ét}}(X,\mathbb Q_\ell)\)
has integral coefficients independent of \(\ell\),
and that every complex conjugate of each eigenvalue
has absolute value \(q^{i/2}\)
\cite[theorem (1.6), p.~276]{Deligne1974WeilI}.
Here the Frobenius action is the one in Deligne's
statement: the pullback induced by the \(q\)-power
geometric endomorphism, equivalently the geometric
Galois Frobenius with his convention.
The finite-field cardinality, the geometric
cohomology, and this action are part of the theorem.
Deligne's later theory of weights is developed in
\cite{Deligne1980WeilII}. These citations acknowledge
his results; they do not assert that the present
chart and Schwartz source already satisfy those
geometric hypotheses.

Exposé-level credits behind this literature are retained in the source
guide. In particular, SGA~7, Exposé~XXI is by Katz, with its integrality
appendix by Deligne \citep{Katz1973SGA7XXI,Deligne1973Integrality}.
Illusie's account records the rédaction and contributor roles in SGA~5
and SGA~4\(\tfrac12\) \citep{IllusieChronologySGA5,Deligne1977SGA4Half}.
These historical credits do not add hypotheses to the theorem above.

The actual local comparison in this programme is
as follows. On a selected zero retain its coordinate
\(z=s-\rho\) and the complete germ
\[
g_\rho(z)=g(\rho+z)=u_\rho(z)z^{m_\rho},
\qquad u_\rho(0)\ne0,\qquad
A_\rho=\mathcal O_\rho/(g_\rho).
\]
The equality of ideals
\((g_\rho)=(z^{m_\rho})\) follows from invertibility
of \(u_\rho\); it supplies Taylor coordinates but
does not remove the unit from any formula.
For a real parameter \(q>0\), define
\[
C_qf(z)=f(z/q)\pmod{g_\rho},\qquad N_\rho f=zf.
\]
This \(q\) is a positive grading parameter, not a
claim of finite-field structure. In particular it
must not be identified solely by its notation with
the polynomial degree denoted \(q\) elsewhere.
The substitution is well-defined because
\[
\frac{g_\rho(z/q)}{g_\rho(z)}
=q^{-m_\rho}\frac{u_\rho(z/q)}{u_\rho(z)}
\]
is a holomorphic unit at zero. Its inverse is
substitution \(z\mapsto qz\).
Thus \(C_pC_q=C_{pq}\), \(C_1=1\), and
\(C_q(z^j)=q^{-j}z^j\) for \(0\le j<m_\rho\).

The original Mellin scaling acts on the same full
block by
\[
U_a|_{A_\rho}
=a^\rho e^{(\log a)N_\rho}
=a^\rho\sum_{j=0}^{m_\rho-1}
\frac{(\log a)^j}{j!}N_\rho^j,\qquad a>0,
\]
because \(\mathcal M(U_aF)(s)=a^s\mathcal MF(s)\).
Direct substitution and multiplication now give
\[
C_qN_\rho C_q^{-1}=q^{-1}N_\rho,\qquad
U_aN_\rho U_a^{-1}=N_\rho,
\]
\[
C_qU_a|_{A_\rho}C_q^{-1}
=a^\rho e^{(\log a)N_\rho/q}
=a^{\rho(1-1/q)}U_{a^{1/q}}|_{A_\rho}.
\]
These are exact morphism identities on the original
local algebra, retaining the complete finite
exponential and its scalar factor.

There is also a precise obstruction to a particular
simultaneous identification. Write
\(A_Z=\bigoplus_{\rho\in Z}A_\rho\),
\(N=\bigoplus N_\rho\), and retain \(U_a\) on that
sum. Suppose a target vector space carries an
invertible \(F\) and an operator \(N_D\) satisfying
\(FN_DF^{-1}=q^{-1}N_D\), with \(q\ne1\).
For every linear map \(T\) intertwining both proposed
operators, namely \(TU_a=FT\) and \(TN=N_DT\),
we have \(F^{-1}T=TU_a^{-1}\), and therefore
\[
q^{-1}N_DT
=FN_DF^{-1}T
=FTNU_a^{-1}
=TU_aNU_a^{-1}
=TN=N_DT.
\]
It follows that \(TN=0\). The unique surviving
factorization is through
\[
A_Z\twoheadrightarrow A_Z/NA_Z
\cong\bigoplus_{\rho\in Z}\mathbb C,
\qquad [f_\rho]\longmapsto(f_\rho(0))_{\rho\in Z}.
\]
The kernel of evaluation in each Taylor block is
exactly \(zA_\rho=N_\rho A_\rho\), proving the
quotient identification. Evaluation intertwines
\(U_a\) with \(\operatorname{diag}(a^\rho)\) and
sends \(N\) to zero, giving a concrete surviving
comparison. An injective simultaneous intertwiner
can exist only when all \(N_\rho=0\), equivalently
all retained orders are one.

This obstruction concerns those two specified
intertwining equations. The full-jet automorphism
\(C_q\) remains available and has already been
constructed above. Its conjugacy relation alone
cannot supply Deligne's absolute weight conclusion:
for any nonzero scalar \(c\), replacing \(C_q\)
by \(cC_q\) leaves the conjugacy with \(N\)
unchanged and multiplies every eigenvalue by \(c\).
On a simple block \(C_q=1\) for every \(q\),
whereas the original scaling still has eigenvalue
\(a^\rho\). The displayed source sections,
support maps, and operator comparisons are the
proved connections. Further geometric and
quantitative constructions remain part of the
research programme, with their original objects
and hypotheses still present.

\section{Reading the human-source record}
\label{app:provenance}
The source record follows a mathematical use, rather than ending at the most
recent project note that contains it. It records the cited work, its human
authors and other documented contributors, the precise result used, the
programme's receiving statement, and the evidence actually inspected. An
earlier source credited by that work receives a further dependency edge when
its contribution is relevant to the argument. Authors of an exposition and
authors of its underlying results retain their respective credits.

The accompanying files distinguish five levels of evidence. A
\emph{primary-text check} means the stated passage of the human source was
read and compared with its application. An \emph{authenticated identity}
establishes the publication and author metadata, without implying that its
theorem was read. An \emph{inherited review} preserves an older source record
with its original scope. A \emph{discovery candidate} is a name, eponym,
formula, or citation awaiting resolution. An \emph{unresolved edge} states
the missing evidence. These labels are attached to specific entries; they
are not interchangeable quality scores for an entire paper.

This distinction matters for the present corpus. A scan of the 692-page predecessor's
15 editable text files covered 223,268 lines. Its bibliography blocks contain
five distinct programme references and no human-authored bibliography entry.
This does not show that human sources were never used: the historical
literature catalogue and recovered reading records contain many such sources.
The companion reconnects those sources to the actual mathematical uses. A
name search by itself establishes neither a dependency nor a theorem's
correctness.

The original acquired public conversation has 5,349 nonempty message records
and 5,350 nodes including its empty root. Its 215 readable conversational
records are a selection; 2,600 tool payloads were redacted in the published
chain. Those missing bodies have not been reconstructed merely because a
reference survives. Other exports and retained mathematical manuscripts are
audited separately, with duplicate bodies identified where possible.
\texttt{COVERAGE.md} records the current boundary and remaining source gaps.
\texttt{ATTRIBUTION\_MAP.md} and \texttt{ATTRIBUTION\_GRAPH.json} connect
uses to source works; \texttt{UNRESOLVED\_ATTRIBUTION.md} gives the exact
remaining checks. The six tau reports and the expanded historical
bibliography preserve the longer recursive source trails.

\subsection{A worked attribution repair: the logarithmic cumulant identity}
The secondary globalization note uses an unnamed standard moment--cumulant
identity. Its exact formula is \citet[\S2.3.4, p.~38, equation~(2.9)]{McCullagh2018Tensor},
whose formal-series derivation applies without positivity assumptions.
A further human reading route is \citet[p.~2]{McCullagh2011SpeedCommentary},
which explains the partition-lattice treatment of \citet{Speed1983Cumulants}
and its connection with the inversion framework of \citet{Rota1964Mobius}.
McCullagh's text was read; the two earlier publication identities were
authenticated separately. Their complete original articles were not newly
read in this check. The following calculation supplies the exact convention
map to the project formula \citep{SplitZeroSecondary2026}.

Let \(F\) be analytic and nonzero near \(t\), and write
\[
 m_n=\frac{F^{(n)}(t)}{F(t)},\qquad
 \kappa_n=(\log F)^{(n)}(t),\qquad
 M(z)=\frac{F(t+z)}{F(t)}=1+\sum_{j\ge1}m_j\frac{z^j}{j!}.
\]
A local analytic logarithm exists. Its positive-order derivatives are
independent of the chosen branch. Expanding its Taylor series gives
\[
 \sum_{n\ge1}\kappa_n\frac{z^n}{n!}
 =\log M(z)
 =\sum_{r\ge1}\frac{(-1)^{r-1}}{r}
       \left(\sum_{j\ge1}m_j\frac{z^j}{j!}\right)^r.
\]
For ordered positive sizes \(j_1+\cdots+j_r=n\), the coefficient
\(n!/(j_1!\cdots j_r!)\) counts ordered partitions of
\(\{1,\ldots,n\}\) into blocks of those sizes. Each unordered partition
with \(r\) blocks occurs \(r!\) times. The factor \(1/r\) therefore leaves
\((r-1)!\), proving
\[
 \kappa_n=\sum_{\pi\in\Pi_n}
        (-1)^{|\pi|-1}(|\pi|-1)!\prod_{B\in\pi}m_{|B|}.
\]
In the note's notation,
\(A_j(t;s)=(s)_jL_t(s+j)/L_t(s)\) and
\(m_j=(-1)^jA_j\). For each partition, the block sizes sum to \(n\).
Consequently its full sign is \((-1)^n\), and the exact receiving identity is
\[
 \partial_t^{,n}\log L_t(s)
 =(-1)^n\sum_{\pi\in\Pi_n}
        (-1)^{|\pi|-1}(|\pi|-1)!\prod_{B\in\pi}A_{|B|}(t;s).
\]
This is a classical combinatorial identity applied in the note's analytic
domain. The complex-valued application requires no probability measure.
The source attribution has been repaired; no priority claim follows from
the project specialization.

\subsection{How to use the companion files}
\texttt{HUMAN\_SOURCES.md} provides the bibliographic and mathematical role
of the resolved sources; \texttt{REFERENCES.bib} retains their reusable
identities. The attribution records include project-source edges and
source-to-source edges with verification levels. The inherited bibliography
is retained as provenance without upgrading all of it to a fresh primary
reading. Historical routes no longer used in the current calculation remain
credited with their historical role.

The public guide excludes private conversation and assignment text.
Source hashes and stable public theorem locations allow a mathematical
comparison without publishing those private records. A source gap remains
visible until its missing identity, passage, or application map is actually
resolved. The mathematical tasks continue while that reconstruction proceeds.

\clearpage
\bibliographystyle{plainnat}
\bibliography{REFERENCES}
\end{document}
